% concatenated from main.tex
\documentclass[11pt,final]{amsart}
\usepackage{amstext, amscd, amsmath, amssymb}
\usepackage{mathtools}
%\usepackage{fullpage} 

% change font style
%\usepackage{palatino}
%\usepackage{times}
%\usepackage{tgbonum}
%\usepackage{charter}
%\usepackage{times}

% check if all items are cited
%\usepackage{refcheck}

%\usepackage[notcite,notref]{showkeys}
%\usepackage{showlabels}
%\usepackage{color}
%\usepackage{hyperref}
%\usepackage{tikz-cd}


% style of text
\numberwithin{equation}{section}
%\setlength{\parskip}{1.05pt}
\footskip=20pt %for numbering of pages

\usepackage[a4paper]{geometry}
\geometry{
    tmargin= 3cm, %3cm, %1.3in, %3.5cm, %2.65cm,
    bmargin= 2.5cm, %2.3cm, %1.1in, %3cm, %2.5cm,
    rmargin= 2.5cm, %2.3cm, %1.1in, %2.75cm, %2.3cm,
    lmargin= 2.5cm %2.3cm %1.1in %2.75cm, %2.3cm
    }
    
%    \usepackage[a4paper]{geometry}
%\geometry{
%    tmargin= 3cm, 
%    bmargin= 3cm, 
%    rmargin= 3cm, 
%    lmargin= 3cm 
%    }


% adjust spaces in bibliography
\let\OLDthebibliography\thebibliography
\renewcommand\thebibliography[1]{
  \OLDthebibliography{#1}
  \setlength{\parskip}{0pt}
  \setlength{\itemsep}{2pt plus 0.5ex}
}

%%%%%%%%%% begin macros %%%%%%%%%%%%%%%
%
%      Cites in bold rather than roman.
\makeatletter
\def\@cite#1#2{{\m@th\upshape\bfseries%
[{#1\if@tempswa{\m@th\upshape\mdseries, #2}\fi}]}}
\makeatother
%
% Proclamation definitions in the most emphatic (plain) style:
\theoremstyle{plain}
\newtheorem{theorem}{Theorem}[section]
\newtheorem{corollary}[theorem]{Corollary}
\newtheorem{proposition}[theorem]{Proposition}
\newtheorem{lemma}[theorem]{Lemma}
% Proclamation definitions in the less emphatic (definition) style:
\theoremstyle{definition}
\newtheorem{definition}[theorem]{Definition}
\newtheorem{example}[theorem]{Example}
\newtheorem{examples}[theorem]{Examples}
\newtheorem{counterexample}[theorem]{Counterexample}
\newtheorem{remark}[theorem]{Remark}
\newtheorem{application}[theorem]{Application}
\newtheorem{question}[theorem]{Question}
\newtheorem*{acknow}{Acknowledgements}
\newtheorem*{details}{Details}
\newtheorem*{open}{Open Access Statement}
% Proclamation definitions in the least emphatic (remark) style:
\theoremstyle{remark}
%\newtheorem{remark}{Remark}
\newtheorem*{notation}{Notation}

%      Proof environment
\renewcommand{\proofname}{\em\textbf{Proof.}}
\renewcommand{\qedsymbol}{{\vrule height5pt width5pt depth1pt}}
% Change first-level `enumerate' numbering from arabic to roman.
%
\renewcommand{\labelenumi}{\textup{(\roman{enumi})} }
% properly align :=
\mathtoolsset{centercolon}
%
%      Capital script letters
  \newcommand{\A}{{\mathcal{A}}}
  \newcommand{\B}{{\mathcal{B}}}
  \newcommand{\C}{{\mathcal{C}}}
  \newcommand{\D}{{\mathcal{D}}}
  \newcommand{\E}{{\mathcal{E}}}
  \newcommand{\F}{{\mathcal{F}}}
  \newcommand{\G}{{\mathcal{G}}}
\renewcommand{\H}{{\mathcal{H}}}
  \newcommand{\I}{{\mathcal{I}}}
  \newcommand{\J}{{\mathcal{J}}}
  \newcommand{\K}{{\mathcal{K}}}
\renewcommand{\L}{{\mathcal{L}}}
  \newcommand{\M}{{\mathcal{M}}}
  \newcommand{\N}{{\mathcal{N}}}
\renewcommand{\O}{{\mathcal{O}}}
\renewcommand{\P}{{\mathcal{P}}}
  \newcommand{\Q}{{\mathcal{Q}}}
  \newcommand{\R}{{\mathcal{R}}}
\renewcommand{\S}{{\mathcal{S}}}
  \newcommand{\T}{{\mathcal{T}}}
  \newcommand{\U}{{\mathcal{U}}}
  \newcommand{\V}{{\mathcal{V}}}
  \newcommand{\W}{{\mathcal{W}}}
  \newcommand{\X}{{\mathcal{X}}}
  \newcommand{\Y}{{\mathcal{Y}}}
  \newcommand{\Z}{{\mathcal{Z}}}
%Greek Letters
\newcommand{\eps}{\varepsilon}
\def\al{\alpha}
\def\be{\beta}
\def\Ga{\Gamma}
\def\ga{\gamma}
\def\De{\Delta}
\def\de{\delta}
\def\ze{\zeta}
\def\io{\iota}
\def\ka{\kappa}
\def\la{\lambda}
\def\La{\Lambda}
\def\om{\omega}
\def\Om{\Omega}
\def\si{\sigma}
\def\Si{\Sigma}
\def\Ups{\Upsilon}
\newcommand\qs{{\rm QS}}
\newcommand\vthe{\vartheta}
\newcommand\vphi{\varphi}
%\def\eps{\epsilon}
\newcommand\wpi{\widetilde{\pi}}
\newcommand\wsi{\widetilde{\sigma}}
%Roman letters for math
\newcommand{\rC}{{\mathrm{C}}}
\newcommand{\rN}{\mathrm{N}}
%Math capital letters
\newcommand{\bB}{\mathbb{B}}
\newcommand{\bC}{\mathbb{C}}
\newcommand{\bD}{\mathbb{D}}
\newcommand{\bF}{\mathbb{F}}
\newcommand{\bK}{\mathbb{K}}
\newcommand{\bN}{\mathbb{N}}
\newcommand{\bT}{\mathbb{T}}
\newcommand{\bZ}{\mathbb{Z}}
\newcommand{\bR}{\mathbb{R}}
%Gothic letters
\newcommand{\fA}{{\mathfrak{A}}}
\newcommand{\fB}{{\mathfrak{B}}}
\newcommand{\fC}{{\mathfrak{C}}}
\newcommand{\fD}{{\mathfrak{D}}}
\newcommand{\fE}{{\mathfrak{E}}}
\newcommand{\fF}{{\mathfrak{F}}}
\newcommand{\fG}{{\mathfrak{G}}}
\newcommand{\fg}{{\mathfrak{g}}}
\newcommand{\fH}{{\mathfrak{H}}}
\newcommand{\fh}{{\mathfrak{h}}}
\newcommand{\fI}{{\mathfrak{I}}}
\newcommand{\fJ}{{\mathfrak{J}}}
\newcommand{\fK}{{\mathfrak{K}}}
\newcommand{\fL}{{\mathfrak{L}}}
\newcommand{\fM}{{\mathfrak{M}}}
\newcommand{\fN}{{\mathfrak{N}}}
\newcommand{\fO}{{\mathfrak{O}}}
\newcommand{\fP}{{\mathfrak{P}}}
\newcommand{\fQ}{{\mathfrak{Q}}}
\newcommand{\fR}{{\mathfrak{R}}}
\newcommand{\fS}{{\mathfrak{S}}}
\newcommand{\fs}{{\mathfrak{s}}}
\newcommand{\fT}{{\mathfrak{T}}}
\newcommand{\fU}{{\mathfrak{U}}}
\newcommand{\fV}{{\mathfrak{V}}}
\newcommand{\fW}{{\mathfrak{W}}}
\newcommand{\fX}{{\mathfrak{X}}}
\newcommand{\fY}{{\mathfrak{Y}}}
\newcommand{\fZ}{{\mathfrak{Z}}}
%      Lower case bold letters
\newcommand{\Ba}{{\mathbf{a}}}
\newcommand{\Bb}{{\mathbf{b}}}
\newcommand{\Bc}{{\mathbf{c}}}
\newcommand{\Bd}{{\mathbf{d}}}
\newcommand{\Be}{{\mathbf{e}}}
\newcommand{\Bf}{{\mathbf{f}}}
\newcommand{\Bg}{{\mathbf{g}}}
\newcommand{\Bh}{{\mathbf{h}}}
\newcommand{\Bi}{{\mathbf{i}}}
\newcommand{\Bj}{{\mathbf{j}}}
\newcommand{\Bk}{{\mathbf{k}}}
\newcommand{\Bl}{{\mathbf{l}}}
\newcommand{\Bm}{{\mathbf{m}}}
\newcommand{\Bn}{{\mathbf{n}}}
\newcommand{\Bo}{{\mathbf{o}}}
\newcommand{\Bp}{{\mathbf{p}}}
\newcommand{\Bq}{{\mathbf{q}}}
\newcommand{\Br}{{\mathbf{r}}}
\newcommand{\Bs}{{\mathbf{s}}}
\newcommand{\Bt}{{\mathbf{t}}}
\newcommand{\Bu}{{\mathbf{u}}}
\newcommand{\Bv}{{\mathbf{v}}}
\newcommand{\Bw}{{\mathbf{w}}}
\newcommand{\Bx}{{\mathbf{x}}}
\newcommand{\By}{{\mathbf{y}}}
\newcommand{\Bz}{{\mathbf{z}}}
%      Text used in equations
\newcommand{\AND}{\text{ and }}
\newcommand{\FOR}{\text{ for }}
\newcommand{\FORAL}{\text{ for all }}
\newcommand{\IF}{\text{ if }}
\newcommand{\IFF}{\text{ if and only if }}
\newcommand{\OR}{\text{ or }}
\newcommand{\foral}{\text{ for all }}
\newcommand{\qand}{\quad\text{and}\quad}
\newcommand{\qif}{\quad\text{if}\quad}
\newcommand{\qiff}{\quad\text{if and only if}\quad}
\newcommand{\qfor}{\quad\text{for}\quad}
\newcommand{\qforal}{\quad\text{for all}\ }
\newcommand{\qforsome}{\quad\text{for some}\ }
\newcommand{\qotherwise}{\quad\text{otherwise}}
%      Useful shortforms
\newcommand{\BH}{{\B(\H)}}
\newcommand{\bsl}{\setminus}
\newcommand{\ca}{\mathrm{C}^*}
\newcommand{\cstar}{\mathrm{C}^*}
\newcommand{\cenv}{\mathrm{C}^*_{\textup{env}}}
\newcommand{\tenv}{\mathcal{T}_{\textup{env}}}
\newcommand{\ltwo}{\ell^2}
\newcommand{\lip}{\langle}
\newcommand{\rip}{\rangle}
\newcommand{\ip}[1]{\lip #1 \rip}
\newcommand{\bip}[1]{\big\lip #1 \big\rip}
\newcommand{\Bip}[1]{\Big\lip #1 \Big\rip}
\newcommand{\ol}{\overline}
\newcommand{\wt}{\widetilde}
\newcommand{\wh}{\widehat}
\newcommand{\sot}{\textsc{sot}}
\newcommand{\wot}{\textsc{wot}}
\newcommand{\wotclos}[1]{\ol{#1}^{\textsc{wot}}}
%      Better math forms
\newcommand{\lh}{\, \lhd \,}
%      Operators
\newcommand{\ad}{\operatorname{ad}}
\newcommand{\Alg}{\operatorname{Alg}}
\newcommand{\ann}{\operatorname{ann}}
\newcommand{\aug}{\operatorname{aug}}
\newcommand{\Aut}{\operatorname{Aut}}
\newcommand{\alg}{\operatorname{alg}}
\newcommand{\ball}{\operatorname{ball}}
\newcommand{\cois}{\operatorname{co-is}}
\newcommand{\cov}{\operatorname{c}}
\newcommand{\core}{\operatorname{ }}
\newcommand{\dirlim}{\varinjlim}
\newcommand{\diag}{\operatorname{diag}}
\newcommand{\Dim}{\operatorname{dim}}
\newcommand{\dist}{\operatorname{dist}}
\newcommand{\dom}{\operatorname{Dom}}
\newcommand{\End}{\operatorname{End}}
\newcommand{\env}{\operatorname{\textup{env}}}
\newcommand{\ev}{\operatorname{ev}}
\newcommand{\fock}{\operatorname{F}}
\newcommand{\GL}{\operatorname{GL}}
\newcommand{\Id}{\operatorname{Id}}
\newcommand{\id}{{\operatorname{id}}}
\newcommand{\im}{{\operatorname{Im}}}
\newcommand{\Ind}{{\operatorname{Ind}}}
\newcommand{\irred}{{\operatorname{irred}}}
\newcommand{\iso}{{\operatorname{is}}}
\newcommand{\LL}{\operatorname{L}}
\newcommand{\Lat}{\operatorname{Lat}}
\newcommand{\mt}{\emptyset}
\newcommand{\Mult}{\operatorname{Mult}}
\newcommand{\nc}{\operatorname{nc}}
\newcommand{\nd}{\operatorname{nd}}
\newcommand{\nil}{\operatorname{nil}}
\newcommand{\PI}{\operatorname{PI}}
\newcommand{\Prim}{\operatorname{Prim}}
\newcommand{\rad}{\operatorname{rad}}
\newcommand{\ran}{\operatorname{ran}}
\newcommand{\rank}{\operatorname{rank}}
%\newcommand{\Ref}{\operatorname{Ref}}
\newcommand{\reg}{{\operatorname{rg}}}
\newcommand{\rep}{\operatorname{rep}}
\newcommand{\scv}{\operatorname{sc}}
\newcommand{\spn}{\operatorname{span}}
\newcommand{\sumoplus}{\operatornamewithlimits{\sum \strut^\oplus}}
\newcommand{\supp}{\operatorname{supp}}
\newcommand{\uni}{{\operatorname{un}}}
%Other
\newcommand{\sca}[1]{\left\langle#1\right\rangle} %\sca{a,b} =<a,b>
\newcommand{\nor}[1]{\left\Vert #1\right\Vert} %\nor{x}=||x||
\newcommand{\bo}[1]{\mathbf{#1}} %bold math symbols
\newcommand{\un}[1]{{\underline{#1}}} %underline math symbols
\newcommand{\ncl}[1]{\left[ #1 \right]^{-\|\cdot\|}} %norm closure linear span
\newcommand{\wcl}[1]{\left[ #1 \right]^{-w^*}} % weak* closure linear span
\newcommand{\quo}[2]{{\raisebox{.1em}{$#1$}\left/ \, \raisebox{-.1em}{$#2$}\right.}} % quotient of objects
%Subshifts
\newcommand{\umu}{\underline\mu}
\newcommand{\unu}{\underline\nu}
\newcommand{\uw}{\underline{w}}
\newcommand{\uq}{\underline{q}}
\newcommand{\umt}{\underline\mt}
\newcommand{\fdn}{(\mathbb{F}_+^d)^N}
\newcommand{\kcnp}{\mathcal{K}_{\textup{CNP}}}

%Do not show subsections in the table of contents
%\addtocontents{toc}{\protect\setcounter{tocdepth}{1}}

%%%%%%%%%%%%%%%%%%%%%%%%%%%%%%%%

%%%%%%%%%%%%%%%%%%%%%%%%%%%%%%%%
\begin{document}

\title[The reduced Hao--Ng isomorphism problem for product systems]{The reduced Hao--Ng isomorphism problem for non-degenerate product systems}

\author[I.A. Paraskevas]{Ioannis Apollon Paraskevas}
\address{Department of Mathematics\\ National and Kapodistrian University of Athens\\ Athens\\ 157 84\\ Greece}
\email{ioparask@math.uoa.gr}

\subjclass[2020]{46L08, 47L55, 46L05, 46K50, 46L55}
\keywords{Product systems, Fock space, Fock covariance, Hao--Ng isomorphism, reduced crossed products of operator algebras}

%%%%%%%%%%%%%%%%%%%%%%%%%%%%%%%%
\begin{abstract}
We prove that the reduced Hao--Ng isomorphism problem has an affirmative answer for a generalised gauge action of a locally compact Hausdorff group on a non-degenerate product system over a unital subsemigroup of a discrete group. In the possible absence of a faithful conditional expectation on the reduced crossed product, we employ the Plancherel weight and integrability of reduced coactions. Using this approach, we establish that the identity representation of the induced product system inside the reduced crossed product of the Fock C*-algebra is Fock covariant.
\end{abstract}

\maketitle

\section{Introduction}

\subsection{Product systems framework}

Product systems over a unital subsemigroup $P$ of a discrete group $G$ provide a framework for modelling irreversible dynamics. They include operator algebras arising from subshifts, graphs and higher-rank graphs, topological graphs, and semigroup actions on C*-algebras. They first appeared in the work of Arveson \cite{Arv89} and were later studied by Dinh \cite{Din91} for discrete subsemigroups of $\bR_+$. Nica \cite{Nic92} considered semigroups with a quasi-lattice order, where the left regular representation has particularly useful properties. The relations coming from the Fock representation, now known as Nica covariance, give a more manageable class of representations. Nica's setting may be viewed as a product system with one-dimensional fibres.

Pimsner \cite{Pim97} developed the theory for $P=\bZ_+$, where a product system is determined by a single C*-correspondence. This construction includes graph C*-algebras, crossed products by $\bZ$, and algebras associated with partial dynamical systems; see for example \cite{Kat03}. Motivated by \cite{Nic92,Pim97} and his work with Raeburn \cite{FR98}, Fowler \cite{Fow02} initiated the study of product systems over quasi-lattice ordered semigroups. Brownlowe, Larsen and Stammeier \cite{BLS18} later considered product systems arising from dynamical systems over right LCM semigroups. Kwa\'{s}niewski and Larsen \cite{KL19a,KL19b} extended this theory to general product systems over right LCM semigroups, which were further studied by Dor-On, Kakariadis, Katsoulis, Laca and Li \cite{DKKLL22}, and Kakariadis, Katsoulis, Laca and Li \cite{KKLL23}. Beyond this setting, Wick ordering and Nica covariance are no longer available. Nevertheless, Cuntz, Deninger and Laca \cite{CDL13} showed that the Fock representation can still be used as the model for covariance.

Coactions by the ambient group $G$ have played a central role in the theory of product systems, since they allow one to use tools from the theory of Fell bundles. Dor-On, Kakariadis, Katsoulis, Laca and Li \cite{DKKLL22} observed that the normal coaction by $G$ on $\T_\la(X)$ gives a canonical Fell bundle, denoted here by $\F\C X$. Its reduced C*-algebra is $\T_\la(X)$, while its full C*-algebra $\T_{\cov}^{\fock}(X)$ need not coincide with the universal Toeplitz algebra $\T(X)$ of all representations of $X$. A representation of $X$ is called Fock covariant when it factors through $\T_{\cov}^{\fock}(X)$. It was shown in \cite{DKKLL22} that, for compactly aligned product systems over right LCM semigroups, Fock covariance agrees with Nica covariance. In this case, $\T_{\cov}^{\fock}(X)$ is the universal Nica covariant C*-algebra $\N\T(X)$. 

Li \cite{Li12}, building on ideas from \cite{CDL13}, introduced relations that record not only the semigroup multiplication but also the principal right ideals and their intersections, known as the constructible right ideals. These relations are modelled by the Fock representation and were studied further by Kakariadis, Katsoulis, Laca and Li \cite{KKLL22}, also in connection with the inverse semigroup models of Norling \cite{Nor14}. Laca and Sehnem \cite{LS22} gave a complete description of the relations that make an isometric representation of a semigroup Fock covariant. In \cite{KP25}, Kakariadis and the author gave a characterisation of injective equivariant Fock covariant representations of a product system without any assumptions on the semigroup or the product system.

A long-standing problem has been whether there exists a boundary quotient of $\T(X)$ for the injective equivariant Fock covariant representations of $X$, and how this quotient relates to the C*-envelope of the tensor algebra $\T_\la(X)^+$, in analogy with the Cuntz--Pimsner algebra $\O_X$ of Katsura \cite{Kat04} for $P=\bZ_+$. Carlsen, Larsen, Sims and Vittadello \cite{CLSV11} proved the existence of such a terminal object in the case of compactly aligned product systems over quasi-lattice ordered pairs $(G,P)$ and under some assumptions that are satisfied by many important cases, such as $(\bZ^d,\bZ_+^d)$. The universal construction for general product systems was given by Sehnem \cite{Seh19}, who introduced the quotient $A\times_X P$. Here $A$ denotes the coefficient algebra of $X$, i.e., the unit fibre $A:= X_{e_P}$. The key feature of $A\times_X P$ is that $A$ embeds in $A\times_X P$, and for a $*$-representation of $A\times_X P$, injectivity on $A$ forces injectivity on the fixed-point algebra of $A\times_X P$. Dor-On and Katsoulis \cite{DK20} then settled the problem for compactly aligned product systems over abelian lattice ordered groups by showing that the terminal object for injective equivariant Fock covariant representations exists and coincides with the C*-envelope of the tensor algebra, and with $A\times_X P$.
In \cite{DKKLL22}, Dor-On, Kakariadis, Katsoulis, Laca, and Li identified $A\times_X P$ with the full C*-algebra of the strong covariant bundle, denoted by $\S\C X$. Furthermore, in the compactly aligned right LCM case, they proved that the terminal object is canonically $*$-isomorphic to both the reduced C*-algebra of this Fell bundle, denoted by $A\times_{X,\la}P$, and the coaction C*-envelope. Note that in the abelian lattice ordered case the full and reduced C*-algebras of $\S\C X$ coincide. Finally, Sehnem \cite{Seh22} proved that completely isometric representations of $\T_\la(X)^+$ automatically admit a conditional expectation. This led to the identification of $A\times_{X,\la}P$ with $\cenv(\T_\la(X)^+)$, which resolved the problem in general and unified all previous results.

\subsection{The reduced Hao--Ng isomorphism problem}
Suppose now that a locally compact Hausdorff group \(\fH\) acts on $\T_\la(X)$ by a generalised gauge action
$\al\colon\fH\to\Aut(\T_\la(X))$, i.e., a point-norm continuous homomorphism such that $\al_\fh(\la_p(X_p))=\la_p(X_p)$ for every $p\in P$ and
$\fh\in\fH$.  There is an induced product system
\[
X\rtimes_{\al,\la}\fH
=
\{X_p\rtimes_{\al,\la}\fH\}_{p\in P}
\]
over $P$, with coefficient algebra
\(A\rtimes_{\al,\la}\fH\), realised in the reduced crossed product $\T_\la(X)\rtimes_{\al,\la}\fH$. The action $\al$ extends to an action $\dot\al$ on the C*-envelope of $\T_\la(X)^+$ which is canonically $*$-isomorphic with
$A\times_{X,\la}P$.  The reduced Hao--Ng isomorphism problem then asks whether the reduced strong covariant functor commutes with the
reduced crossed product functor, namely whether there is a
canonical $*$-isomorphism
\[
(A\rtimes_{\al,\la}\fH)
\times_{X\rtimes_{\al,\la}\fH,\la}P
\ \simeq\
(A\times_{X,\la}P)
\rtimes_{\dot\al,\la}\fH.
\]

We begin by discussing the problem in the case $P=\mathbb Z_+$, where the product system $X$ is determined by a single C*-correspondence, and the preceding question reduces to the original problem for Cuntz--Pimsner algebras as introduced by Hao and Ng \cite{HN08}.
Hao and Ng \cite{HN08} established the corresponding $*$-isomorphism for actions of
amenable locally compact Hausdorff groups. As applications of their results, they recover previous results on Hilbert bimodules by Abadie \cite{Aba10}, and on graph C*-algebras by Kumjian and Pask \cite{KP99}.
Furthermore, Katsoulis \cite{Kat17} connected the Hao--Ng isomorphism problem with the work of Echterhoff, Kaliszewski, Quigg and Raeburn \cite{EKQR06} on imprimitivity theorems for C*-dynamical systems.

The amenability assumption
was subsequently weakened in several directions. In particular, B\'edos, Kaliszewski,
Quigg and Robertson \cite{BKQR15} resolved the problem for actions of exact discrete groups. In a series of works initiated by Katsoulis and Ramsey \cite{KR19} and followed by \cite{Kat17, KR21}, the authors developed a very effective non-selfadjoint approach to the reduced Hao--Ng isomorphism problem. In \cite{KR21} it was proved that in the case of $P=\bZ_+$ and for any locally compact Hausdorff group $\fH$ there exists a canonical $*$-isomorphism
\[
\T_\la(X \rtimes_{\al, \la} \fH) {\simeq} \T_\la(X) \rtimes_{\al, \la} \fH,
\]
that restricts to a completely isometric isomorphism
\[
\T_\la(X \rtimes_{\al, \la} \fH)^+ {\simeq} \T_\la(X)^+ \rtimes_{\al, \la} \fH. 
\]
This was previously shown when $\fH$ is discrete in \cite{Kat17} and when $\fH$ is locally compact Hausdorff and abelian in \cite{KR19}.
By employing the identification of the C*-envelope of the tensor algebra of a C*-correspondence obtained by Katsoulis and Kribs in \cite{KK06}, the problem is reduced in the case $P=\bZ_+$ to determining whether there exists a canonical $*$-isomorphism
\begin{equation} \label{eq:cenvcp}
\cenv(\T_\la(X)^+ \rtimes_{\al, \la} \fH)\stackrel{?}{\simeq} \cenv(\T_\la(X)^+) \rtimes_{\dot{\al}, \la} \fH.    
\end{equation}

This isomorphism \eqref{eq:cenvcp} was established in \cite{Kat17} for any operator algebra that admits a contractive approximate unit (in the place of $\T_\la(X)^+$) when $\fH$ is discrete. In \cite{KP25} Kakariadis and the author removed the assumption of the contractive approximate unit. In \cite{KR19} it was established for operator algebras that admit a contractive approximate unit and $\fH$ locally compact Hausdorff and abelian, and in \cite{KR21} it was established for a locally compact Hausdorff group $\fH$ and hyperrigid operator algebras (in the sense of Arveson \cite{Arv11}) that admit a contractive approximate unit.  In a recent work,
Dor-On and Thompson \cite{DT26} proved \eqref{eq:cenvcp} for any locally compact Hausdorff group and any operator algebra that admits a contractive approximate unit, hence they resolved the reduced Hao--Ng isomorphism problem for all
non-degenerate product systems over $P=\bZ_+$.

We now move the discussion to product systems over more general semigroups. The reduced Hao--Ng isomorphism problem was resolved for a compactly aligned product system over an abelian lattice order $(G,P)$, by Dor-On and Katsoulis \cite{DK20} when $\fH$ is a discrete group, and subsequently by Katsoulis \cite{Kat20} when $\fH$ is a locally compact Hausdorff and abelian group.
It was then resolved by Dor-On, Kakariadis, Katsoulis, Laca and Li \cite{DKKLL22} for a compactly aligned product system over a right LCM semigroup $P$ with $\fH$ discrete. Kakariadis and the author \cite{KP25} resolved the problem for any product system over a unital subsemigroup $P$ of a discrete group $G$ and $\fH$ discrete. We note that in \cite{KP25} the product system $X$ is not assumed to be non-degenerate. The non-degeneracy of $X$ is usually assumed in order to obtain a contractive approximate unit for the tensor algebra $\T_\la(X)^+$ from $A$.

The use of the C*-envelope toolkit for the reduced Hao–Ng isomorphism problem stands as an essential example of the interaction between selfadjoint and non-selfadjoint operator algebras. This approach---pioneered by Katsoulis and Ramsey \cite{KR19}, and followed by many researchers in numerous subsequent works \cite{DKKLL22, DK20, KP25, Kat17, Kat20, KR21} including the present paper—demonstrates how co-universal representations in the
category of non-selfadjoint operator algebras yield a canonical $*$-isomorphism, providing a solution to a selfadjoint problem.

\subsection{Approach and main results}

The purpose of this work is to settle the reduced Hao--Ng isomorphism problem for a non-degenerate product system $X$ in full generality, i.e., for any unital subsemigroup $P$ of a discrete group $G$ and any locally compact Hausdorff group $\fH$.

The main approach in \cite{DKKLL22, DK20, Kat20} has been to use the independence condition for right LCM semigroups and compact alignment of the product system to show Nica covariance of the identity representation 
\begin{equation}\label{eq:ident}
X \rtimes_{\al, \la} \fH \hookrightarrow \T_\la(X) \rtimes_{\al, \la} \fH,
\end{equation}
in order to obtain a canonical $*$-isomorphism
\begin{equation}\label{eq:fccp}
\T_\la(X \rtimes_{\al, \la} \fH) \simeq \T_\la(X) \rtimes_{\al, \la} \fH.
\end{equation}
From there it follows that
\begin{equation}\label{eq:tencp}
\T_\la(X \rtimes_{\al, \la} \fH)^+ \simeq \T_\la(X)^+ \rtimes_{\al, \la} \fH,
\end{equation}
and then applying the C*-envelope machinery concludes the proof. Moving beyond the right LCM condition and in the absence of independence and compact alignment it is unclear whether the $*$-isomorphism \eqref{eq:fccp} holds. However, the completely isometric isomorphism \eqref{eq:tencp} is enough in order to conclude the proof.
In \cite{KP25}, Kakariadis and the author proved that \eqref{eq:tencp}
holds for any product system $X$ when $\fH$ is discrete, by showing that the identity representation
\eqref{eq:ident} is Fock covariant. The main obstruction in passing from the discrete case
to the locally compact Hausdorff case is the proof that the identity representation is
Fock covariant. In particular, in \cite{KP25} the discreteness of $\fH$ is used
to obtain a faithful conditional expectation on the reduced crossed product
$\T_\la(X)\rtimes_{\al,\la}\fH$. Together with the characterisation of
injective equivariant Fock covariant representations obtained in \cite{KP25}, this
yields Fock covariance of the identity representation.

In order to overcome this difficulty for locally compact Hausdorff groups, we work with the averaging map associated to a slice map of the Plancherel weight of the reduced group C*-algebra $\ca_\la(\fH)$ and the reduced dual coaction by $\fH$ on the reduced crossed product $\T_\la(X)\rtimes_{\al,\la}\fH$; see Subsection \ref{Ss:Aver}. This averaging map was introduced and studied by Buss and Meyer, and Buss, in \cite{BM09,Bus10} in their study of integrability of coactions, and we use it in order to replace the faithful conditional expectation that one has in the discrete case.  After multiplying by a suitable compactly supported approximate unit obtained by non-degeneracy of $X$, arbitrary elements of the reduced crossed product $\T_\la(X)\rtimes_{\al,\la}\fH$ are localised inside the square-integrable domain of the coaction, where the averaging map is well-defined.

In particular, in Lemma~\ref{L:local1}, Lemma~\ref{L:invar} and Lemma~\ref{L:local2}, we prove the essential properties of the averaging map and develop this localisation approach, that are needed to prove our main result.

\medskip

\noindent
\textbf{Theorem~A.}
(Theorem~\ref{T:idisfock})
{\it 
Let $X$ be a non-degenerate product system over a unital subsemigroup $P$ of a discrete group
$G$, and let $\al$ be a generalised gauge action of a locally compact Hausdorff group
$\fH$ on $\T_\la(X)$.  The identity representation
\[
\io^\rtimes \colon X\rtimes_{\al,\la}\fH
\to
\T_\la(X)\rtimes_{\al,\la}\fH
\]
is an injective Fock covariant representation that admits a normal coaction by $G$.
}

\medskip

With this at hand, applying the results of Dor-On and Thompson \cite{DT26} and Sehnem \cite{Seh22} gives the asserted canonical $*$-isomorphism. In particular, we obtain the positive resolution of the reduced Hao--Ng isomorphism problem for non-degenerate product systems in full generality.

\medskip

\noindent
\textbf{Theorem~B.}
(Theorem~\ref{T:HaoNg})
{\it 
Let $P$ be a unital subsemigroup of a discrete group $G$ and let $X$ be a non-degenerate product system over $P$.
Let $\al$ be a generalised gauge action of a locally compact Hausdorff group $\fH$ on $\T_\la(X)$.
Then the identity representation $\io^\rtimes \colon X\rtimes_{\al,\la}\fH \to \T_{\la}(X) \rtimes_{\al,\la} \fH$ induces a completely isometric isomorphism
\[
\T_\la(X \rtimes_{\al, \la} \fH)^+ \simeq \T_\la(X)^+ \rtimes_{\al, \la} \fH.
\]
Consequently, the reduced Hao--Ng isomorphism problem has an affirmative answer, i.e.,
\[
(A \rtimes_{\al, \la} \fH) \times_{X \rtimes_{\al, \la} \fH,\la} P
\simeq
(A \times_{X, \la} P) \rtimes_{\dot \al, \la} \fH,
\]
by a canonical $*$-isomorphism, where $\dot{\al}$ is the induced action of $\fH$ on $A \times_{X,\la} P$.
}

\subsection{Contents}
The paper is organised as follows.  In Section 2 we recall the required
material on operator algebras, coactions, crossed products and Fell bundles, while we also fix notation. Section 3 is
devoted to weights, slice maps and integrable elements of coactions, and we establish some properties of the averaging map that are needed in the proof of the main result.
In Section 4 we recall the required facts on product systems, and employ our localisation approach to prove that the identity representation of the induced product system inside the reduced crossed product of the Fock C*-algebra is Fock covariant. We then apply the C*-envelope machinery to resolve the reduced Hao--Ng isomorphism problem for non-degenerate product systems.

\subsection{Acknowledgements}
The author is grateful to Evgenios Kakariadis for fruitful discussions, as well as for his careful reading and comments and suggestions on a draft of this manuscript. The author is also grateful to Michalis Anoussis, Alexandros Chatzinikolaou and Aristides Katavolos for their comments and suggestions on a draft of this manuscript.
The author acknowledges that this research work was supported by the Hellenic Foundation for Research and Innovation (HFRI) under the 5th Call for HFRI PhD Fellowships (Fellowship Number: 19145).
This material is based upon work supported by the Swedish Research Council under grant no. 2021-06594 while the author was in residence at Institut Mittag-Leffler in Djursholm, Sweden during the first semester of 2026.
The author would like to thank the Institut Mittag-Leffler and the organisers of the program ``Operator Algebras and Quantum Information'' for the hospitality.

\section{Preliminaries}

We write $\otimes$ for the minimal tensor product between C*-algebras.
A map between algebras with given sets of generators is called canonical if it preserves generators of the same index.
If $\X$ is a subset of a normed linear space $\Y$, then we write $[\X]$ for the closed linear span of $\X$ in $\Y$. The dual of a Banach space $\X$ is denoted by $\X^d$. To avoid any possible confusion, we always denote discrete groups by $G$ and locally compact Hausdorff groups by $\fH$. When we consider a unital subsemigroup $P$ of a group $G$ we always mean that $e_P=e_G$. We write $\C_+$ for the positive cone of a C*-algebra $\C$.

%%%%%%%%%%%%%%%%%%%%%%%%%%%%%%%%
\subsection{Operator algebras}
%%%%%%%%%%%%%%%%%%%%%%%%%%%%%%%%

For a comprehensive treatment and full details regarding the general theory of non-selfadjoint operator algebras, we refer the reader to \cite{BL04, Pau02}. 

By an operator algebra $\A$, we mean a norm-closed subalgebra of some $\B(H)$.
Any operator algebra possesses a C*-cover, i.e., a pair $(\C,j)$, where $\C$ is a C*-algebra and $j \colon \A \to \C$ is a completely isometric homomorphism satisfying $\C = \ca(j(\A))$. 
The \emph{C*-envelope} $(\cenv(\A),\io)$ is the co-universal C*-cover of $\A$. This means that, given any other C*-cover $(\C,j)$, there exists a unique $*$-epimorphism $\Phi \colon \C \to \cenv(\A)$ such that $\Phi \circ j = \io$. 
Hamana established the existence of the C*-envelope in \cite{Ham79}.

Many of the C*-algebras we will need in this work are going to be non-unital and we will often consider their multiplier algebras. If $\C$ is a C*-algebra we denote by $\M(\C)$ the \emph{multiplier algebra} of $\C$. If we consider $\C$ canonically as a Hilbert $\C$-module, then $\C\simeq \K(\C)$ and $\M(\C)\simeq \L(\C)$ where $\K(\C)$ denotes the generalised compact operators and $\L(\C)$ denotes the adjointable operators; see \cite[Sections~1 and~2]{Lan95}. 
In addition, if $\pi \colon \C \to \B(H)$ is a faithful non-degenerate $*$-representation, then $\pi$ extends to an injective unital $*$-homomorphism $\M(\C) \to \B(H)$ such that
\[
\M(\C)\simeq \{ a\in \B(H) : a \pi(c), \pi(c)a \in \pi(\C) \foral c\in \C\},
\]
see \cite[Proposition 3.12.3]{Ped79}. In particular, $\C$ is an ideal in the unital C*-algebra $\M(\C)$ that is \emph{essential}, i.e., if $x\in \M(\C)$ and $x c=0$ for all $c\in \C$, then $x=0$.  We note that if $\C$ is unital, then $\C=\M(\C)$.
The \emph{strict topology} on $\M(\C)$ is the topology induced by the seminorms $x \to \|xc\|$ and $x \to \|c x\|$ for $c\in \C$.

Let $\C$ and $\D$ be C*-algebras. A $*$-homomorphism $\Phi \colon \C \to \M(\D)$ is called \emph{non-degenerate} if $[\Phi(\C) \D]=\D$. This is also equivalent to the net $(\Phi(e_i))_i$ converging strictly to $1_{\M(\D)}$ for any approximate unit $(e_i)_i$ of $\C$. In this case, $\Phi$ extends to a strictly continuous unital $*$-homomorphism $\ol{\Phi} \colon \M(\C)\to \M(\D)$. If moreover, $\Phi$ is injective, then $\ol{\Phi}$ is also injective. For a proof of these facts
see \cite[Proposition 2.1 and Proposition 2.5]{Lan95}. We will consider the C*-algebras $\C\otimes \M(\D)$, $\M(\C)\otimes \D$ and $\M(\C)\otimes \M(\D)$ inside $\M(\C\otimes \D)$ without further mention. This is justified due to injectivity of the minimal tensor product of C*-algebras, and the existence of an injective unital $*$-homomorphism $\M(\C)\otimes \M(\D) \hookrightarrow\M(\C\otimes \D)$; see \cite[Pages 35--37]{Lan95}.

%%%%%%%%%%%%%%%%%%%%%%%%%%%%%%%%
\subsection{Coactions and Fell bundles}
%%%%%%%%%%%%%%%%%%%%%%%%%%%%%%%%

Let $\fH$ be a locally compact Hausdorff group with unit $e_{\fH}$. We denote the left Haar measure on $\fH$ by $\mu$ and we write $\De$ for the modular function of $\fH$.
We denote by $\ca(\fH)$ and $\ca_\la(\fH)$ the full group C*-algebra and the reduced group C*-algebra of $\fH$, respectively. 
We write $u_{\fH} \colon \fH \to \U(\M(\ca(\fH)))$ for the canonical embedding of $\fH$ as unitary multipliers and $\la_{\fH} \colon \ca(\fH) \to \ca_\la(\fH)$ for the left regular representation of $\fH$. Similarly, we can consider $\fH$ canonically inside $\U(\M(\ca_\la(\fH)))$, and by an abuse of notation, for a group element $\fh\in \fH$ we will write $\la_{\fH}(\fh)$ for the resulting unitary multiplier inside $\U(\M(\ca_\la(\fH)))$.
We recall that for a non-degenerate $*$-homomorphism $\Phi$ defined on a C*-algebra, we write $\ol{\Phi}$ for its strictly continuous extension on the multiplier algebra.

By the universal property of $\ca(\fH)$, there exists an injective non-degenerate $*$-homomorphism
\[
\de_\fH \colon \ca(\fH) \to \M(\ca(\fH) \otimes \ca(\fH)) 
\text{ such that }
(\ol{\de_\fH\otimes \id_{\ca(\fH)}})\circ \de_\fH=(\ol{ \id_{\ca(\fH)} \otimes \de_\fH})\circ \de_\fH.
\]
This map is the integrated form of the strictly continuous homomorphism $\fh \mapsto u_{\fH}(\fh) \otimes u_{\fH}(\fh)$.

Analogously, there is an injective non-degenerate $*$-homomorphism 
\[
\de^r_\fH \colon \ca_\la(\fH) \to \M(\ca_\la(\fH) \otimes \ca_\la(\fH))
\text{ such that }
(\ol{\de^r_\fH\otimes \id_{\ca_\la(\fH)}})\circ \de^r_\fH=(\ol{\id_{\ca_\la(\fH)} \otimes \de^r_\fH})\circ \de^r_\fH.
\]
The canonical extension of this map to the multiplier algebra $\M(\ca_\la(\fH))$ satisfies the relation
\[
\ol{\de^r_\fH}(\la_{\fH}(\fh)) = \la_{\fH}(\fh)\otimes \la_{\fH}(\fh) \foral \fh \in \fH;
\]
for further details see \cite[Pages 141--142]{Kat84}.

We say that a C*-algebra $\C$ \emph{admits a (full) coaction $\de$ by $\fH$} if there exists an injective non-degenerate $*$-homomorphism $\de \colon \C \to \M(\C \otimes \ca(\fH))$ such that 
\[
\de(\C)(1_{\M(\C)}\otimes \ca(\fH)) \subseteq \C\otimes \ca(\fH),
\]
and satisfying
\[
(\ol{\de \otimes \id_{\ca(\fH)}}) \circ \de = (\ol{\id_{\C} \otimes\de_\fH}) \circ \de.
\]
Similarly, we say that a C*-algebra $\C$ \emph{admits a reduced coaction $\de^r$ by $\fH$} if there exists an injective non-degenerate $*$-homomorphism $\de^r \colon \C \to \M(\C \otimes \ca_\la(\fH))$ such that 
\[
\de^r(\C)(1_{\M(\C)}\otimes \ca_\la(\fH)) \subseteq \C\otimes \ca_\la(\fH),
\]
and satisfying
\[
(\ol{\de^r \otimes \id_{\ca_\la(\fH)}}) \circ \de^r = (\ol{\id_{\C} \otimes \de^r_\fH}) \circ \de^r.
\]
For a  coaction $\de$ by $\fH$ on a C*-algebra $\C$ we write
\[
\M^{\rm fix}_\de:=\{x \in \M(\C) : \ol{\de}(x)=x \otimes 1_{\M(\ca(\fH))}\},
\]
for the \emph{fixed-point algebra}.
Similarly, for a reduced coaction $\de^r$ by $\fH$ we write
\[
\M^{\rm fix}_{\de^r}:=\{x \in \M(\C) : \ol{\de^r}(x)=x \otimes 1_{\M(\ca_\la(\fH))}\}.
\]

If $G$ is a discrete group, then $\ca(G)$ and $\ca_\la(G)$ are unital and therefore we can remove non-degeneracy and the inclusion conditions from the above definitions in order to obtain a coaction or a reduced coaction, see for example \cite[Remark 3.2]{DKKLL22}.
In this case, for a coaction $\de$ by $G$, the coaction identity is equivalent to the spectral spaces spanning a dense subset of $\C$, where the \emph{$g$-spectral space} is defined as 
\[
[\C]_g := \{c \in \C \mid \de(c) = c \otimes u_G(g)\} \foral g \in G;
\]
see the proof of \cite[Proposition 2.6]{Ng96}. Note that the fixed-point algebra of $\de$ is exactly the $e_G$-spectral space $[\C]_{e_G}$. In particular, the collection $\{[\C]_g\}_{g\in G}$ is a Fell bundle over $G$ (which we will define below).
If, in addition, the map $(\id_{\C}\otimes \la_G) \circ \de$ is injective, then $\de$ is called \emph{normal}.
It follows that $\de$ is normal if and only if $(\id_{\C}\otimes \la_G) \circ \de$ is a reduced coaction by $G$ on $\C$.

For more details on Fell bundles see \cite{Exe17, FD88}.
A \emph{Fell bundle} over a locally compact Hausdorff group $\fH$ is a continuous Banach bundle $\B = \{\B_{\fh}\}_{\fh \in \fH}$ such that:
\begin{enumerate}
\item there are continuous bilinear and associative \emph{multiplication maps} from $\B_{\fh} \times \B_{\fg}$ to $\B_{\fh\fg}$ such that $\|b_{\fh} b_{\fg} \| \leq \|b_{\fh}\| \cdot \|b_{\fg}\|$;
\item there are continuous conjugate linear \emph{involution maps} from $\B_{\fh}$ to $\B_{\fh^{-1}}$ such that $(b_{\fh}^*)^* = b_{\fh}$ and $\|b_{\fh}^*\| = \|b_{\fh}\|$;
\item $(b_{\fh} b_{\fg})^* = (b_{\fg})^* b_{\fh}^*$;
\item $\|b_{\fh}^* b_{\fh}\| = \| b_{\fh}\|^2$;
\item $b_{\fh}^* b_{\fh} \geq 0$ in $\B_{e_\fH}$.
\end{enumerate}

We denote by $C_c(\B)$ the space of compactly supported continuous sections of $\B$ and make it a $*$-algebra by setting
\[
(F \ast F')(\fg):= \int_{\fH} F(\fh) F'(\fh^{-1}\fg) \, d\mu(\fh) \foral F,F'\in C_c(\B),
\]
and
\[
F^*(\fg):=\De(\fg)^{-1} F(\fg^{-1})^* \foral F\in C_c(\B).
\]
The \emph{full C*-algebra} of $\B$, denoted by $\ca(\B)$, is the enveloping C*-algebra of $C_c(\B)$, i.e., the completion of $C_c(\B)$ with respect to the norm 
\[
\|F\|:=\sup\{ p(F): p \text{ is a C*-seminorm on } C_c(\B)\} \foral F\in C_c(\B).
\]

We denote by $L^2(\B)$ the Hilbert $\B_{e_\fH}$-module defined as the completion of $C_c(\B)$ with respect to the $\B_{e_\fH}$-valued inner product 
\[
\sca{F,F'}:=\int_\fH F(\fh)^* F'(\fh) \, d\mu(\fh) \foral F,F' \in C_c(\B).
\]
By universality of $\ca(\B)$ we obtain a $*$-representation
\[
\la_\B \colon \ca(\B) \to \L(L^2(\B)) \text{ such that } \la_\B(F)(F')=F \ast F'\foral F,F' \in C_c(\B),
\]
which we call the \emph{left regular representation of $\B$}. The \emph{reduced cross-sectional C*-algebra} of $\B$ is then defined to be the C*-algebra $\ca_\la(\B):=\la_\B(\ca(\B))$. We note that the $*$-algebra $C_c(\B)$ is dense in $\ca(\B)$. Moreover, $\la_{\B}$ is injective on $C_c(\B)$ and $\la_\B(C_c(\B))$ is dense in $\ca_\la(\B)$, see \cite[Section~2]{EN02} for a proof of these facts.

In \cite[Proposition 2.10]{EN02} it is shown that if $\B$ is a Fell bundle over a locally compact Hausdorff group $\fH$, then the C*-algebra $\ca_{\la}(\B)$ admits a reduced coaction \[
\de_\B^r \colon \ca_\la(\B) \to \M(\ca_\la(\B)\otimes \ca_\la(\fH)).
\]
The non-degenerate $*$-homomorphism $\de_\B^r$ extends to an injective strictly continuous unital $*$-homomorphism $\ol{\de_\B^r} \colon \M(\ca_\la(\B)) \to \M(\ca_\la(\B)\otimes \ca_\la(\fH))$,  and satisfies
\[
\ol{\de_\B^r}(\ol{\la_{\B}}(b_{\fh}))
=
\ol{\la_{\B}}(b_{\fh})\otimes\la_{\fH}(\fh)
\foral  b_{\fh}\in\B_{\fh} \text{ and all } \fh\in \fH,
\]
where $b_{\fh}$ is considered inside $\M(\ca(\B))$ by \cite[Lemma~1.1]{EN02}.

\subsection{Reduced crossed products}\label{Ss:Cross}
We briefly discuss the properties of the reduced crossed product by a locally compact Hausdorff group that we are going to need later on. For more details see \cite[Chapter 7]{Wil07} and \cite[Appendix A]{EKQR06}.

Once again, let $\fH$ be a locally compact Hausdorff group with unit $e_{\fH}$, the left Haar measure $\mu$, and $\De$ the modular function of $\fH$.
Let $\C$ be a C*-algebra and $\al \colon \fH \to \Aut(\C)$ be a point-norm continuous action. We call the triple $(\C,\fH,\al)$ a \emph{C*-dynamical system}. The space 
$C_c(\fH,\C)$ of continuous $\C$-valued functions on $\fH$ with compact support becomes a $*$-algebra  by setting
\[
(f \ast h)(\fg')
:=
\int_{\fH}f(\fg)\al_{\fg}(h(\fg^{-1}\fg')) \, d\mu(\fg)
\foral f,h\in C_c(\fH,\C),
\]
and
\[
f^*(\fg)
:=
\De(\fg)^{-1}\al_{\fg}(f(\fg^{-1})^*) \foral f\in C_c(\fH,\C).
\]

We denote the \emph{reduced crossed product} by $\C\rtimes_{\al,\la}\fH$.
See \cite[Definition~7.7 and Lemma~7.8]{Wil07} or
\cite[Appendix~A, Definition~A.13 and Remark~A.15]{EKQR06} for its
construction.  By \cite[Appendix~A, Definition~A.13]{EKQR06}, there is a
canonical covariant pair $(i_{\C}^r,i_{\fH}^r)$, where
\[
i_{\C}^r\colon\C\to\M(\C\rtimes_{\al,\la}\fH)
\]
is a non-degenerate injective $*$-homomorphism and
\[
i_{\fH}^r\colon\fH\to\U(\M(\C\rtimes_{\al,\la}\fH))
\]
is a strictly continuous injective group homomorphism into the unitary group $\U(\M(\C\rtimes_{\al,\la}\fH))$ of $\M(\C\rtimes_{\al,\la}\fH)$, satisfying
\[
i_{\fH}^r(\fh)i_{\C}^r(c)i_{\fH}^r(\fh)^*
=
i_{\C}^r(\al_{\fh}(c))
\foral c\in\C \text{ and } \fh\in\fH.
\]

The integrated form 
\[
f
\longmapsto
(i_{\C}^r\times i_{\fH}^r)(f)
:=
\int_{\fH} i^r_{\C}(f(\fh)) i_\fH^r(\fh) \, d\mu(\fh)
\]
of $(i_{\C}^r,i_{\fH}^r)$ is injective on $C_c(\fH,\C)$, and hence it identifies $C_c(\fH,\C)$ with a dense
$*$-subalgebra of $\C\rtimes_{\al,\la}\fH$.
We suppress the use of the above map and regard $C_c(\fH,\C)$ as a dense $*$-subalgebra of $\C\rtimes_{\al,\la}\fH$.
Moreover, note that since $i_{\C}^r$ is non-degenerate, it extends uniquely to a strictly continuous injective unital $*$-homomorphism
\[
\ol{i_{\C}^r}\colon\M(\C)\to\M(\C\rtimes_{\al,\la}\fH).
\]

We use the standard identification of the reduced crossed product of a C*-dynamical system, with the reduced cross-sectional C*-algebra of the corresponding semidirect-product Fell bundle.  Namely, if $(\C,\fH,\al)$ is  a C*-dynamical system, then the \emph{semidirect-product Fell bundle} associated with $(\C,\fH,\al)$ is defined as
\[
\B^\al:=\{\B^\al_{\fh}\}_{\fh\in\fH} \text{ where } \B^\al_{\fh}:=\C\times\{\fh\} \foral \fh \in \fH,
\]
with the operations
\[
(b_{\fh},\fh)(c_{\fg},\fg):=(b_{\fh}\al_{\fh}(c_{\fg}),\fh\fg)
\qand
(b_{\fh},\fh)^*:=(\al_{\fh^{-1}}(b_{\fh}^*),\fh^{-1}).
\]
There exists a canonical $*$-isomorphism 
\[
\Phi \colon \ca_\la(\B^\al) \to \C\rtimes_{\al,\la}\fH.
\]
See for example \cite[Section~2]{EN02}.

Recall that we have a reduced coaction $\de_{\B^\al}^r \colon \ca_\la(\B^\al) \to \M(\ca_\la(\B^\al) \otimes \ca_\la(\fH))$. Moreover, we also have a reduced coaction $\de^{\rtimes} \colon \C\rtimes_{\al,\la}\fH \to \M(\C\rtimes_{\al,\la}\fH \otimes \ca_\la(\fH))$ called the \emph{reduced dual coaction} that satisfies
\[
\ol{\de^\rtimes}(i_\C^r(c))= i_\C^r(c) \otimes 1_{\M(\ca_\la(\fH))} 
\qand 
\ol{\de^\rtimes}(i_\fH^r(\fh))= i_\fH^r(\fh) \otimes \la_\fH(\fh);
\]
see \cite[Pages 254--256]{Lan79}.
In particular, we have $i_\C^r(\C) \subseteq \M^{\rm fix}_{\de^{\rtimes}}$, and
the $*$-isomorphism $\Phi$ intertwines the coactions $\de^\rtimes$ and $\de_{\B^\al}^r$, that is, 
\[
\de^\rtimes \circ \Phi = \ol{\Phi \otimes \id_{\ca_\la(\fH)}} \circ \de_{\B^\al}^r.
\]
The proof follows by an adaptation of the proof of \cite[Proposition 3.5]{KMQW10}.

We close this section with a short discussion on crossed products of operator algebras, following the work of Katsoulis and Ramsey \cite{KR19}.
We recall only the parts that will be needed for the reduced Hao--Ng isomorphism problem.

Suppose that $\fH$ is a locally compact Hausdorff group acting point-norm continuously on an operator algebra $\A$ via completely isometric automorphisms.
This action admits an extension to an action $\dot{\al}$ on the C*-envelope of $\A$, which allows us to construct the reduced crossed product $\cenv(\A) \rtimes_{\dot\al, \la} \fH$.
By identifying $\A$ with its copy inside $\cenv(\A)$, the \emph{reduced crossed product} $\A \rtimes_{\al, \la} \fH$ is defined to be the norm-closed subalgebra of $\cenv(\A) \rtimes_{\dot{\al}, \la} \fH$ generated by the compactly supported continuous functions $C_c(\fH,\A)$, see \cite[Definition 3.17]{KR19}.
A central question in this theory is whether this inclusion induces a canonical $*$-isomorphism with the C*-envelope of $\A \rtimes_{\al, \la} \fH$, or in other words, whether
\begin{equation*}
\cenv(\A \rtimes_{\al, \la} \fH) \stackrel{?}{\simeq} \cenv(\A) \rtimes_{\dot\al, \la} \fH.
\end{equation*}
Recently, Dor-On and Thompson \cite[Theorem~4.4]{DT26} gave an affirmative answer when $\A$ admits a contractive approximate unit, thus generalising the results obtained in \cite[Theorem~2.5]{Kat17}, \cite[Theorem~3.23]{KR19}, and the first part of \cite[Theorem~3.6]{KR21}.

\section{Weights and the averaging map}

\subsection{Weights on C*-algebras}

We briefly review the theory of weights on C*-algebras. The standard reference for lower semi-continuous weights is \cite{Com68}. We follow the notation of \cite{KV99}.

Let $\C$ be a C*-algebra.  A \emph{weight} on
$\C$ is a map $\om \colon \C_+\to[0,\infty]$ which is additive and positively
homogeneous.  
We set
\[
\M_\om^+ := \{c\in \C_+:\om(c)<\infty\},
\]
for the \emph{positive integrable domain} of $\om$.
From the additivity of $\om$ and since its range lies in $[0,\infty]$ we get that if $c,d\in \C_+$ are such that $c\leq d$, then $\om(c)\leq \om(d)$. Hence $\M_\om^+$ is a \emph{hereditary cone} in $\C_+$, i.e., if $c\in \C_+$ and $d\in \M_\om^+$ with $c\leq d$, then $c\in \M_\om^+$.
Furthermore, we set
\[
\N_\om := \{c\in \C:\om(c^*c)<\infty\}
\qand 
\M_\om:=\spn \M_{\om}^+,
\]
for the \emph{square-integrable domain} and the \emph{integrable domain} of $\om$, respectively.
It follows that $\N_\om$ is a left
ideal in $\C$ and that $\M_\om$ is a $*$-subalgebra of $\C$. Moreover, we have 
\[
\M_\om=\spn\N_\om^* \N_\om
\qand 
\M_\om\cap \C_+=\M_\om^+.
\]
It follows that $\om|_{\M_{\om}^+}$ has a unique
linear extension to $\M_\om$, which we still denote by $\om$. See \cite[Lemma~5.1.2]{Ped79} for a proof of these facts.

A weight $\om$ is called \emph{faithful} if whenever $\om(c)=0$ for some $c\in \C_+$ then $c=0$.  It is called \emph{lower semi-continuous} if
$\{c\in \C_+:\om(c)\leq r\}$ is norm-closed for every $r\geq 0$ and \emph{densely defined} if $\N_\om$ is dense in $\C$. A weight that is non-zero, lower semi-continuous, and densely defined is called \emph{proper}.

For a proper weight $\om$ we set 
\[
\F_\om:= \{ \rho \in \C^d: \rho \text{ is positive and } \rho(c) \leq \om(c) \foral c\in \C_+\},
\]
and
\[
\G_{\om}:= \{\be \rho : \be \in (0,1) \text{ and } \rho \in \F_{\om}\}.
\]
Then $\G_{\om}$ is a directed set with respect to the natural order defined by positivity of  functionals on $\C^d$, see \cite[Proposition 3.5]{Kus97A}.
In particular, in \cite[Proposition 1.7]{Com68} it is shown that 
\[
\om(c)= \sup\{ \rho(c) : \rho\in \F_\om\}=\lim_{\rho \in \G_\om} \rho(c) \foral c\in \C_+.
\]
The weight $\om$ has a canonical extension to the positive cone of the multiplier algebra $\M(\C)_+$, see  \cite[Discussion prior to Lemma~2.4]{Kus97A}.
For $x\in \M(\C)_+$ we set 
\[
\ol{\om}(x):= \sup\{ \ol{\rho} (x) : \rho\in \F_\om\} \in [0,\infty],
\]
where $\ol{\rho}$ denotes the unique strictly continuous extension of $\rho$ to $\M(\C)$. It follows that
\[
\ol{\om}(x)= \sup\{ \ol{\rho} (x) : \rho\in \F_\om\}= \sup\{ \ol{\rho} (x) : \rho\in \G_\om\}= \lim_{\rho \in \G_\om} \ol{\rho}(x) \foral x\in \M(\C)_+.
\]
In particular, the map $\ol{\om}$ is additive and positively homogeneous, and hence a weight.

\begin{proposition}\label{P:faithext}
Let $\om$ be a faithful proper weight on a C*-algebra $\C$. Then $\ol{\om}$ is faithful.
\end{proposition}

\begin{proof}
Consider $x \in \M(\C)_+$ such that $\ol{\om}(x)=0$. Fix $c\in \C$.  We have
\[
 0\leq x^{1/2}c^* c x^{1/2}\leq \|c\|^2 x,
\]
and therefore
\[
 0\leq\om(x^{1/2}c^* cx^{1/2})
 \leq \|c\|^2 \ol{\om}(x)=0.
\]
Faithfulness of $\om$ gives $x^{1/2} c^*c x^{1/2}=0$. In particular, we get $x^{1/2} c^*=0$. Since $c$ was arbitrary we obtain $x^{1/2}=0$ and hence $x=0$, as required.
\end{proof}

\subsection{The Plancherel weight}

For the remainder of this section, we fix a locally compact Hausdorff group $\fH$ with unit $e_{\fH}$ and left Haar measure $\mu$, and we let $\De$ denote its modular function.

Our main focus will be on the Plancherel weight defined on the positive cone of the reduced C*-algebra $\ca_{\la}(\fH)_+$; see \cite[Section~VII.2 and VII.3]{Tak03}  for more details.
The \emph{von Neumann algebra of the group} $\fH$ is defined as the bicommutant of the reduced group C*-algebra, i.e., $L(\fH):= \ca_{\la}(\fH)''$.  A vector $\xi\in L^2(\fH)$ is called \emph{left bounded} if the map
\[
C_c(\fH) \to L^2(\fH); \quad f \mapsto \xi \ast f
\]
extends to a bounded operator on $L^2(\fH)$, which we denote by $\la_{\fH}(\xi)$. Using for example \cite[Proposition VII.3.1]{Tak03} we get that if $\xi \in L^2(\fH)$ is left bounded, then $\la_{\fH}(\xi)\in L(\fH)$. We note that the use of the symbol $\la_{\fH}$ is justified by the fact that for $f\in C_c(\fH)\subseteq L^2(\fH)$ we have that $\la_{\fH}(f)$ is exactly the image of $f$ under the left regular representation $\la_{\fH} \colon \ca(\fH) \to \ca_\la(\fH)$.

The \emph{Plancherel weight} $\wt{\om}_{\fH} \colon L(\fH)_+ \to [0,\infty]$ is then defined 
by setting
\[
\wt{\om}_{\fH}(x):= 
\begin{cases}
\sca{\xi,\xi}_{L^2(\fH)} & \text{ if } x^{1/2}=\la_{\fH}(\xi) \text{ for some left bounded } \xi \in L^2(\fH),\\
\infty & \text{ otherwise.}
\end{cases}
\]
It follows that $\wt{\om}_{\fH}$ is a faithful weight and
\[
\N_{\wt{\om}_{\fH}}=\{ \la_{\fH}(\xi): \xi \in L^2(\fH) \text{ is left bounded}\}
\]
and
\[
\wt{\om}_{\fH}(\la_{\fH}(\xi)^*\la_{\fH}(\eta))
=\sca{\xi,\eta}_{L^2(\fH)} \text{ for left bounded }  \xi,\eta\in L^2(\fH).
\]
See \cite[Theorem~VII.2.5]{Tak03} and \cite[Definition VII.3.2]{Tak03} for a proof of the fact that $\wt{\om}_{\fH}$ is well-defined and a proof of the above assertions.

The \emph{C*-algebraic Plancherel weight}
$\om_{\fH}$ on $\ca_{\la}(\fH)$ is defined to be the restriction of $\wt{\om}_{\fH}$ to
$\ca_{\la}(\fH)_+$.
By definition we have 
\[
\N_{\om_{\fH}}=\ca_\la(\fH)\cap \N_{\wt{\om}_{\fH}}= \{ \la_{\fH}(\xi): \xi \in L^2(\fH) \text{ is left bounded and }\la_{\fH}(\xi)\in \ca_\la(\fH)\}.
\]

The weight $\om_{\fH}$ is densely defined since $\la_{\fH}(C_c(\fH)) \subseteq \N_{\om_{\fH}}$. Indeed, by  \cite[Proposition 2.40]{Fol16} we get that every $f\in C_c(\fH)$ is left bounded since
\[
\|f\ast h\|_2 \leq \|f\|_1 \|h\|_2 \foral h \in C_c(\fH).
\]
Moreover, we have $\la_{\fH}(f)\in \ca_\la(\fH)$ and hence $\la_{\fH}(f)\in\N_{\om_{\fH}}$. 
By \cite[Theorem~VII.1.11]{Tak03} we get that $\om_{\fH}$ is lower semi-continuous, and hence the weight $\om_\fH$ is proper.
We denote the extension of $\om_{\fH}$ to the multiplier algebra $\M(\ca_\la(\fH))$ by $\ol{\om}_{\fH}$. We then have
\[
\N_{\ol{\om}_{\fH}}=\N_{\wt{\om}_{\fH}} \cap \M(\ca_\la(\fH))=\{ \la_{\fH}(\xi): \xi \in L^2(\fH) \text{ is left bounded and }\la_{\fH}(\xi)\in \M(\ca_\la(\fH))\}.
\]
Note here that we used the fact that $\M(\ca_\la(\fH))\subseteq L(\fH)$, which follows by the Bicommutant Theorem and \cite[Proposition 3.12.3]{Ped79}, and hence $\ol{\om}_{\fH}=\wt{\om}_{\fH}|_{\M(\ca_\la(\fH))_+}$.

\subsection{Slice maps of the Plancherel weight}
Let $\C$ be a C*-algebra. For any $\rho \in \G_{\om_{\fH}}$ we consider the slice map 
\[
\id_\C \otimes \rho \colon \C \otimes \ca_\la(\fH) \to \C \text{ such that } c\otimes x \mapsto c \rho(x).
\]
This slice map admits a strictly continuous extension  
\[
\ol{\id_\C \otimes \rho} \colon \M(\C \otimes \ca_\la(\fH)) \to \M(\C),
\]
see for example \cite[Lemma~A.30]{EKQR06}.
We define the \emph{slice map of the Plancherel weight}
\[
(\id_{\C}\otimes \om_{\fH})(x):= \text{strict-}\lim_{\rho \in \G_{\om_{\fH}}} (\ol{\id_\C \otimes \rho})(x) \foral x\in \M_{\id_{\C}\otimes \om_{\fH}}^+,
\]
where $\M_{\id_{\C}\otimes \om_{\fH}}^+$ denotes the \emph{positive integrable domain} of $\id_{\C}\otimes \om_{\fH}$ defined as
\[
\M_{\id_{\C}\otimes \om_{\fH}}^+:= \{ x\in \M(\C \otimes \ca_\la(\fH))_+: \left((\ol{\id_\C \otimes \rho})(x)\right)_{\rho \in \G_{\om_{\fH}}} \text{ converges strictly in } \M(\C)\}.
\]
Moreover, we define the \emph{square-integrable domain} of $\id_{\C}\otimes \om_{\fH}$ to be
\[
\N_{\id_{\C}\otimes \om_{\fH}}:=\{ x\in \M(\C\otimes \ca_\la(\fH)): x^*x \in \M_{\id_{\C}\otimes \om_{\fH}}^+\},
\]
and the \emph{integrable domain} of $\id_{\C}\otimes \om_{\fH}$ to be
\[
\M_{\id_{\C}\otimes \om_{\fH}}:=\spn \M_{\id_{\C}\otimes \om_{\fH}}^+.
\]

It is proven in \cite[Result 3.6 and Notation 3.7]{KV99}  that the set $\M_{\id_{\C}\otimes \om_{\fH}}^+$ is a hereditary cone in $\M(\C\otimes \ca_\la(\fH))_+$, that $\M_{\id_{\C}\otimes \om_{\fH}}$ is a $*$-subalgebra of $\M(\C\otimes \ca_\la(\fH))$, and that $\N_{\id_{\C}\otimes \om_{\fH}}$ is a left ideal in $\M(\C\otimes \ca_\la(\fH))$. Moreover, we have
\[
\M_{\id_{\C}\otimes \om_{\fH}}=\spn (\N_{\id_{\C}\otimes \om_{\fH}})^* \N_{\id_{\C}\otimes \om_{\fH}}
\qand
\M_{\id_{\C}\otimes \om_{\fH}}^+= \M_{\id_{\C}\otimes \om_{\fH}} \cap \M(\C\otimes \ca_\la(\fH))_+,
\]
and the map $(\id_{\C}\otimes \om_{\fH})|_{\M_{\id_\C\otimes \om_{\fH}}^+}$ has a unique linear extension to $\M_{\id_{\C}\otimes \om_{\fH}}$, which we denote by the same symbol.
Also, for $x\in \M_{\id_{\C}\otimes \om_{\fH}}^+$ we have that $(\id_{\C}\otimes \om_{\fH})(x)$ is a positive element.

The following proposition is just an adaptation of the proof of \cite[Lemma~A.30]{EKQR06}. We include a proof for completeness.

\begin{proposition}\label{P:posslic}\cite[Lemma~A.30]{EKQR06}
Let $\C$ and $\D$ be C*-algebras and let $x$ be in $\M(\C\otimes \D)$. If $(\ol{\psi \otimes \id_{\D}})(x)=0$ for every positive functional $\psi \colon \C \to \bC$, then $x=0$.
\end{proposition}

\begin{proof}
Consider $\C \subseteq \B(H)$ and $\D \subseteq \B(K)$ acting non-degenerately. Then we may assume that $\M(\C\otimes \D) \subseteq \B(H\otimes K)$.
For $\ze\in H$, define the positive functional
\[
\psi_\ze \colon \C \to \bC \text{ such that } \psi_\ze(c)=\sca{\ze,c \, \ze}
\foral c\in \C.
\]
By checking first on simple tensors and taking norm and then strict limits we get
\[
\sca{\ze\otimes k_1,y(\ze\otimes k_2)}
=
\sca{k_1,(\ol{\psi_{\ze} \otimes \id_{\D}})(y) k_2} \foral y\in \M(\C\otimes\D) \text{ and } \ze \in H \text{ and } k_1,k_2 \in K.
\]
In particular, we have
\[
\sca{\ze\otimes k_1,x (\ze\otimes k_2)}
=
\sca{k_1,(\ol{\psi_{\ze} \otimes \id_{\D}})(x) k_2}
=0
\foral \ze\in H \text{ and } k_1,k_2\in K.
\]
A polarisation argument then yields that $x=0$, as required.
\end{proof}

We will need the following proposition.

\begin{proposition}\label{P:faithslice}
If $x\in \M_{\id_{\C}\otimes \om_{\fH}}^+$ satisfies $(\id_{\C}\otimes \om_{\fH})(x)=0$, then $x=0$. 
\end{proposition}

\begin{proof}
Suppose that $\psi \colon \C \to \bC$ is a positive functional and denote by $\ol{\psi}$ its strictly continuous extension to $\M(\C)$. Consider now the positive map $\psi \otimes \id_{\ca_\la(\fH)}$ and its strictly continuous extension $\ol{\psi \otimes \id_{\ca_\la(\fH)}}$ to $\M(\C\otimes \ca_\la(\fH))$ which is also positive. Combining with \cite[Proposition 3.9]{KV99} we get  
\[
(\ol{\psi \otimes \id_{\ca_\la(\fH)}}) (x) \in \M_{\ol{\om}_{\fH}} \cap \M(\ca_\la(\fH))_+=\M_{\ol{\om}_{\fH}}^+,
\]
and also 
\[
\ol{\om}_{\fH}((\ol{\psi\otimes \id_{\ca_\la(\fH)}}) (x))=\ol{\psi}((\id_{\C}\otimes \om_{\fH})(x))=0.
\]
By Proposition \ref{P:faithext} we get $(\ol{\psi \otimes \id_{\ca_\la(\fH)}}) (x)=0$ and hence Proposition \ref{P:posslic} yields that $x=0$, as required.
\end{proof}

\subsection{The averaging map}\label{Ss:Aver}

Let $\C$ be a C*-algebra that admits a reduced coaction $\de^r$ by $\fH$ and consider its unique strictly continuous extension to an injective unital $*$-homomorphism 
\[
\ol{\de^r}\colon \M(\C)\to \M(\C\otimes \ca_\la(\fH)).
\]

Following \cite{BM09, Bus10} we say that a positive element $x\in \M(\C)_+$ is \emph{integrable} if $\ol{\de^r}(x)\in \M_{\id_{\C}\otimes \om_{\fH}}^+$. We define \emph{the positive integrable domain of $\de^r$} as
\[
\M_{\de^r}^+:=\{ x\in \M(\C)_+: \ol{\de^r}(x) \in \M_{\id_{\C}\otimes \om_{\fH}}^+\},
\]
the \emph{square-integrable domain of $\de^r$} as
\[
\N_{\de^r}:=\{ x\in \M(\C): x^*x \in \M_{\de^r}^+\},
\]
and the \emph{integrable domain of $\de^r$} as
\[
\M_{\de^r}:= \spn \M_{\de^r}^+.
\]

The following lemma is remarked in \cite[Section~3.2]{BM09}.
We include a proof for completeness.

\begin{lemma}\cite[Section~3.2]{BM09}\label{L:properties}
With the aforementioned notation $\M_{\de^r}^+$ is a hereditary cone in $\M(\C)_+$, and $\M_{\de^r}$ is a $*$-subalgebra of $\M(\C)$. Moreover, $\N_{\de^r}$ is a left ideal in $\M(\C)$ and  
\[
\M_{\de^r}= \spn \N_{\de^r}^*\N_{\de^r} \qand
\M_{\de^r}^+=\M_{\de^r}\cap \M(\C)_+.
\]
\end{lemma}

\begin{proof}
First note that since $\ol{\de^r}$ is a $*$-homomorphism we get that  $\M_{\de^r}^+$ is a hereditary cone in $\M(\C)_+$ as $\M_{\id_{\C}\otimes \om_{\fH}}^+$ is a hereditary cone in $\M(\C\otimes \ca_\la(\fH))_+$. Similarly, $\N_{\de^r}$ is a left ideal in $\M(\C)$ as $\N_{\id_{\C}\otimes \om_{\fH}}$ is a left ideal in $\M(\C\otimes \ca_\la(\fH))$. Using the inequality
\[
(x+ \be y)^* (x+\be y) \leq 2 x^* x+ 2 |\be|^2 y^*y \text{ for } x,y\in \N_{\de^r} \text{ and }\be \in \bC,
\]
combined with the fact that $\M_{\de^r}^+$ is a hereditary cone, yields that $\N_{\de^r}$ is a linear subspace. For $x \in \M_{\de^r}^+$ we have $x^{1/2}\in \N_{\de^r}$, and hence 
\[
x=(x^{1/2})^* x^{1/2} \in \spn\N_{\de^r}^* \N_{\de^r},
\]
showing that $\M_{\de^r} \subseteq \spn \N_{\de^r}^*\N_{\de^r}$. For the converse inclusion pick $x,y\in \N_{\de^r}$. By polarisation we get
\[
x^*y=\frac{1}{4} \sum_{k=0}^3 i^{k}(x+i^k y)^*(x+i^k y) \in \spn \M_{\de^r}^+=\M_{\de^r},
\]
and thus
\[
\M_{\de^r}=\spn \N_{\de^r}^* \N_{\de^r}.
\]
We have that $\M_{\de^r}$ is obviously $*$-closed. Using the fact that $\N_{\de^r}$ is a left ideal yields that $\M_{\de^r}$ is closed under multiplication, and hence it is a $*$-subalgebra of $\M(\C)$. 
We have already shown that if $x \in \M_{\de^r}^+$, then $x\in \spn \N_{\de^r}^*\N_{\de^r}= \M_{\de^r}$ and hence $\M_{\de^r}^+\subseteq \M_{\de^r}\cap \M(\C)_+$. On the other hand, if $x \in \M_{\de^r}\cap \M(\C)_+$, then 
\[
\ol{\de^r}(x) \in \M(\C\otimes \ca_\la(\fH))_+ \qand \ol{\de^r}(x)
\in
\M_{\id_{\C}\otimes \om_{\fH}}. 
\]
In particular, we have
\[
\ol{\de^r}(x)
\in
\M_{\id_{\C}\otimes \om_{\fH}}
\cap
\M(\C\otimes \ca_\la(\fH))_+
=
\M_{\id_{\C}\otimes \om_{\fH}}^+,
\]
and hence $x\in \M_{\de^r}^+$.
\end{proof}

Following \cite[Section~3.2]{Bus10} the \emph{averaging map (with respect to $\de^r$)} is defined to be the map
\[
\E_{\de^r}\colon \M_{\de^r} \to \M(\C); \quad x \mapsto (\id_{\C}\otimes \om_{\fH})(\ol{\de^r}(x)).
\]
Note that $\E_{\de^r}$ is linear as $\id_\C\otimes\om_{\fH}$ is linear on $\M_{\id_\C\otimes\om_{\fH}}$.
Now fix a C*-dynamical system $(\C,\fH,\al)$, and recall the reduced dual coaction $\de^\rtimes$ by $\fH$ on the reduced crossed product $\C\rtimes_{\al,\la} \fH$ defined in Section~\ref{Ss:Cross}.
In \cite[Corollary 3.12]{Bus10} it is shown that
\begin{equation}\label{Eq:Buss}
C_c(\fH,\C)\subseteq\N_{\de^{\rtimes}},
\end{equation}
and that for $f,h \in C_c(\fH, \C)$ we have
\begin{equation}\label{Eq:avercomp}
\E_{\de^{\rtimes}}(f^* \ast h)= i_\C^r\left((f^*\ast h)(e_{\fH})\right)= i_\C^r\left(
\int_{\fH}
\al_{\fg^{-1}}\big(f(\fg)^*h(\fg)\big) d\mu(\fg)\right).
\end{equation}
In particular, we obtain
\[
\E_{\de^{\rtimes}}(f^* \ast h) \in i_\C^r(\C) \qand \E_{\de^{\rtimes}}(f^* \ast f)\in i_\C^r(\C_+).
\]

We note that \cite[Corollary 3.12]{Bus10} is stated in the more general context of reduced cross-sectional C*-algebras of Fell bundles. Here we use the semidirect-product Fell bundle picture and the canonical $*$-isomorphism $\Phi \colon \ca_\la(\B^\al) \to \C\rtimes_{\al,\la}\fH$ that intertwines the reduced coactions in order to translate \cite[Corollary 3.12]{Bus10} into the reduced crossed product setting. See the discussion in Section~\ref{Ss:Cross}.

\begin{lemma}\label{L:Bim}
With the aforementioned notation, the following hold:
\begin{enumerate}
\item If $c,d\in \C$ and $x\in\M_{\de^{\rtimes}}$, then $i_\C^r (c)x i_\C^r (d)\in \M_{\de^{\rtimes}}$, and 
\[
\E_{\de^{\rtimes}}(i_\C^r (c)x i_\C^r (d))=i_\C^r(c) \E_{\de^{\rtimes}}(x) i_\C^r(d).
\]
\item If $x\in\N_{\de^{\rtimes}}$ and $\E_{\de^{\rtimes}}(x^*x)=0$, then $x=0$.
\end{enumerate}
\end{lemma}

\begin{proof}
To prove item (i) recall that $i_\C^r(\C) \subseteq \M^{\rm fix}_{\de^{\rtimes}}$ and apply \cite[Proposition 3.15 (ii) and (iii)]{Bus10}.
For item (ii) note that if $\E_{\de^{\rtimes}}(x^*x)=0$, then
\[
(\id_\C\otimes\om_{\fH})(\ol{\de^{\rtimes}}(x^* x))= \E_{\de^{\rtimes}}(x^*x) =0.
\]
Proposition \ref{P:faithslice} and injectivity of $\ol{\de^{\rtimes}}$ imply that $x^*x=0$, and hence $x=0$.
\end{proof}

\section{The main result}

\subsection{Product systems and their representations}

In his seminal paper, Fowler \cite{Fow02} formalised product systems, motivated by the works of Nica \cite{Nic92}, Pimsner \cite{Pim97}, as well as his joint work with Raeburn \cite{FR98}. We will work in a more general framework of concrete product systems, as considered in \cite{KP25}. For a more detailed description of concrete product systems, alongside a justification of how they generalise the product systems originally formulated in \cite{Fow02}, see \cite[Section~2]{KP25}.

Let $P$ be a unital subsemigroup of a discrete group $G$ with unit $e_P=e_G$.
We say that a family $X = \{X_p\}_{p \in P}$ of closed operator spaces in a common $\B(H)$ is a \emph{(concrete) product system} if the following are satisfied:
\begin{enumerate}
\item $A := X_{e_P}$ is a C*-algebra (called the \emph{coefficient algebra of $X$});
\item $X_p \cdot X_q \subseteq X_{pq}$ for all $p, q \in P$;
\item $X_p^* \cdot X_{pq} \subseteq X_q$ for all $p, q \in P$,
\end{enumerate}
where $\cdot$ means operator multiplication in $\B(H)$.
Uniqueness of $q \in P$ in item (iii) follows since $P$ is left-cancellative.
If, in addition, $[A\cdot X_p]=X_p$ for every $p\in P$, then $X$ is called \emph{non-degenerate}.

The Fock representation of a product system was first considered in \cite{Fow02}. See also \cite{KP25} for the definition of the Fock representation in the context of concrete product systems. Consider the \emph{Fock space} $\F X := \sum^\oplus_{r \in P} X_r$, as a right Hilbert $A$-module.
For every $\xi_p \in X_p$, we define the \emph{left creation operator}
\[
\la_p(\xi_p)\colon \F X \to \F X \text{ such that } \la_p(\xi_p) \eta_r = \xi_p \cdot \eta_r \text{ for } \eta_r \in X_r \subseteq \F X.
\]
It follows that each  $\la_p(\xi_p)$ is adjointable and 
\[
\la_p(\xi_p)^* \eta_s =
\begin{cases}
\xi_p^* \cdot \eta_s & \text{if } s \in pP, \\
0 & \text{if } s \notin pP
\end{cases}
\quad \text{ for } \eta_s \in X_s \subseteq \F X.
\]
We write $\la \colon X \to \L(\F X)$ for this map, and we refer to $\la=\{\la_p\}_{p\in P}$ as the \emph{Fock representation of $X$}.
We write $\T_\la(X)$ for the \emph{Fock C*-algebra of $X$} defined as $\ca(\la_p(X_p) \mid p \in P)$.

A \emph{representation $t = \{t_p\}_{p \in P}$ of $X$} is a family of linear maps $t_p \colon X_p \to \B(K)$ such that:
\begin{enumerate}
\item $t_{e_P}$ is a $*$-representation of $A := X_{e_P}$;
\item $t_p(\xi_p) t_q(\xi_q) = t_{pq}(\xi_p \xi_q)$ for all $\xi_p \in X_p$ and $\xi_q \in X_q$;
\item $t_p(\xi_p)^* t_{pq}(\xi_{pq}) = t_q(\xi_p^* \xi_{pq})$ for all $\xi_p \in X_p$ and $\xi_{pq} \in X_{pq}$.
\end{enumerate}
A representation $t$ is called \emph{injective} if $t_{e_P}$ is injective on $A$.
It follows that the Fock representation is an injective representation of $X$.
The \emph{Toeplitz algebra $\T(X)$ of $X$} is the universal C*-algebra generated by the \emph{universal representation $\wh{t}$ of $X$} with respect to representations of $X$.
Let $t \colon X \to \B(K)$ be a representation of $X$, then there is an induced $*$-representation
\[
t_\ast \colon \T(X) \to \B(K); \quad \wh{t}_p(\xi_p)\mapsto t_p(\xi_p),
\]
and we write
\[
\ca(t) := t_\ast(\T(X)) = \ca(t_p(X_p) \mid p \in P).
\]

Following \cite{Li12}, for a set $Z \subseteq P$ and $p \in P$ we write
\[
pZ:=\{px \mid x \in Z\}
\qand
p^{-1}Z := \{y \in P \mid py \in Z\}.
\]
We write $\J$ for the smallest family of right ideals of $P$ containing $P$ and $\mt$ that is closed under left multiplication, taking pre-images under left multiplication (in the sense above) and finite intersections.
The right ideals of $P$ in $\J$ are called \emph{constructible}.
By the proof of \cite[Lemma~3.3]{Li12} the set of constructible ideals is automatically closed under finite intersections and hence
\[
\J = \left\{ q_n^{-1} p_n \dots q_1^{-1} p_1 P \mid n \in \bN; p_i, q_i \in P \right\} \cup \{\mt\}.
\]

A representation $t = \{t_p\}_{p \in P}$ of $X$ is called \emph{equivariant} if there exists a $*$-homomorphism
\[
\de \colon \ca(t) \to \ca(t) \otimes \ca(G) ; \quad  t_p(\xi_p) \mapsto t_p(\xi_p) \otimes u_{G}(p).
\]
It follows that $\de$ is injective with a left inverse given by the map $\id \otimes \chi$ where $\chi \colon \ca(G)\to \bC$ is the character obtained from the trivial unitary representation of $G$ by universality.
Moreover, $\de$ satisfies the coaction identity and hence $\ca(t)$ admits a coaction by $G$.
For simplicity, we say that \emph{$t$ admits a coaction by $G$} if such a $\de$ exists. 
In this case, the $g$-spectral subspace $[\ca(t)]_g$, for $g \in G$, is the closed linear span of the elements
\begin{align*}
 t_{p_1}(X_{p_1})^* t_{q_1}(X_{q_1}) \cdots t_{p_n}(X_{p_n})^*t_{q_n} (X_{q_n})
\text{ such that } 
p_1^{-1} q_1 \cdots p_n^{-1} q_n = g \text{ and } n \in \bN.
\end{align*}
A proof for this realisation of the $g$-spectral space can be found in \cite[Lemma~2.2]{Seh19} for non-degenerate product systems. Similar arguments give the conclusion in the general case. In \cite[Lemma~2.2]{Seh19} it is shown that $\wh{t}$ admits a coaction by $G$ and in \cite[Proposition~4.1]{DKKLL22} it is shown that $\la$ admits a normal coaction by $G$.

Let $t$ be a representation of $X$.
For $\bo{x} \in \J$ we define the \emph{$\bo{K}$-core on $\bo{x}$} of $\ca(t)$ to be the closed linear span of the spaces
\[
t_{p_1}(X_{p_1})^*  t_{q_1}(X_{q_1}) \cdots t_{p_n}(X_{p_n})^* 
t_{q_n} (X_{q_n})
\]
for any $n\in\bN$ and $p_1, q_1, \dots, p_n, q_n \in P$ that satisfy
\[
p_1^{-1} q_1 \cdots p_n^{-1} q_n = e_P
\text{ and }
q_n^{-1} p_n \dots q_1^{-1} p_1 P = \bo{x}.
\]
We write $\bo{K}_{\bo{x}, t_\ast}$ for this closed space.
We note that we do not require that $\bo{K}_{\mt, t_{\ast}} = (0)$.
Each $\bo{K}_{\bo{x}, t_\ast}$ is a C*-subalgebra of $\ca(t)$ satisfying $\bo{K}_{\bo{x}, t_\ast}=t_\ast(\bo{K}_{\bo{x}, \wh{t}_\ast})$, and moreover, it follows that  
\[
\bo{K}_{\bo{x}, t_\ast} \cdot \bo{K}_{\bo{y}, t_\ast} \subseteq \bo{K}_{\bo{x}\cap \bo{y}, t_\ast} \foral \bo{x}, \bo{y} \in \J.
\]

\subsection{Fock covariant and strong covariant representations}
We now review the key elements on Fock covariant and strong covariant representations. We recall only the parts that will be needed for the reduced Hao--Ng isomorphism problem.
For more details see \cite{DKKLL22, KP25, Seh19, Seh22}.

Let $P$ be a unital subsemigroup of a discrete group $G$ and let $X$ be a product system over $P$. In \cite[Proposition~4.1]{DKKLL22} it is shown that $\la$ admits a normal coaction and hence we obtain a Fell bundle
\[
\F\C X := \Big\{ [\T_\la(X)]_g \Big\}_{g \in G},
\]
consisting of the spectral subspaces of this coaction.
We call $\F\C X$ the \emph{Fock covariant} bundle of $X$. In particular, we have 
\[
\T_\la(X) \simeq \ca_\la(\F\C X),
\]
and we also set 
\[
\T_{\cov}^{\fock}(X):=\ca(\F\C X).
\]
Moreover, $\T_{\cov}^{\fock}(X)$ is a quotient of $\T(X)$ that inherits a coaction by $G$.
A \emph{Fock covariant representation $t$ of $X$} is a representation of $X$ for which the induced $*$-representation $t_\ast$ factors through $\T_{\cov}^{\fock}(X)$. It is easy to see that the Fock representation $\la$ is a Fock covariant representation.

In \cite[Theorem~3.2]{KP25} Kakariadis and the author proved the following characterisation for injective equivariant Fock covariant representations. Even though it is not needed for our purposes in this work, we note that an inspection of the arguments in the proof of \cite[Theorem~3.2]{KP25} shows that equivariance is neither used nor required.

\begin{theorem}\cite[Theorem~3.2]{KP25}\label{T:Fockcov}
Let $X$ be a product system over a unital subsemigroup $P$ of a discrete group $G$. An injective (equivariant) representation $t$ of $X$ is Fock covariant if and only if $t$ satisfies the following conditions:
\begin{enumerate}
\item $\bo{K}_{\mt, t_\ast} = (0)$.
\item For any $\cap$-closed $\F=\{\bo{x}_1,\dots,\bo{x}_n\} \subseteq \J$ such that $\bigcup_{i=1}^n \bo{x}_i \neq \mt$, and any $b_{\bo{x}_i}\in \bo{K}_{\bo{x}_i,\wh{t}_\ast}$, with $i=1,\dots,n$, the following property holds:
\begin{center}
if $\sum\limits_{i: r\in \bo{x}_i}t_\ast(b_{\bo{x}_i})t_r(X_r)=(0)$ for all $r\in \bigcup_{i=1}^n \bo{x}_i$, then $\sum\limits_{i=1} ^n t_\ast(b_{\bo{x}_i})=0$.    
\end{center}
\end{enumerate}
\end{theorem}

We now move the discussion to strong covariant representations.
We write $A \times_X P$ for the \emph{strong covariant C*-algebra of $X$}.
It is shown in \cite{Seh19} that $A \times_X P$ inherits a coaction by $G$.
The \emph{strong covariant} bundle of $X$ was first considered in \cite{DKKLL22}, and it is defined as the Fell bundle
\[
\S\C X := \Big\{ [A \times_X P]_g \Big\}_{g \in G},
\]
consisting of the spectral subspaces of the coaction by $G$ on $A \times_X P$.
By universality there exists a canonical $*$-isomorphism $A \times_X P\simeq \ca(\S\C X)$.
We write $A \times_{X, \la} P$ for the reduced C*-algebra of $\S\C X$.
A \emph{strong covariant representation $t$ of $X$} is a representation of $X$ for which the induced $*$-representation $t_\ast$ factors through $A \times_X P$. In \cite[Theorem~3.10]{Seh19} it is shown that if $t$ is a strong covariant representation then $t$ is injective if and only if $t_\ast$ is injective on $[A \times_X P]_{e_P}$.

For a product system $X$ we write $\T_\la(X)^+$ for the \emph{tensor algebra of $X$}, i.e., for the norm-closed subalgebra generated by $\{\la_p(X_p)\}_{p \in P}$ inside $\T_\la(X)$.
In \cite[Theorem~5.1]{Seh22} Sehnem proved that the C*-envelope $\cenv(\T_\la(X)^+)$ admits a normal coaction by $G$ and thus there exists a canonical $*$-isomorphism 
\[
\cenv(\T_\la(X)^+) \simeq A \times_{X, \la} P.
\]

\subsection{The reduced Hao--Ng isomorphism problem}

Let $P$ be a unital subsemigroup of a discrete group $G$, and let $X$ be a
product system over $P$.  A \emph{generalised gauge action} of a locally compact Hausdorff group $\fH$ on $\T_\la(X)$ is a point-norm continuous homomorphism
$\al \colon \fH\to\Aut(\T_\la(X))$ satisfying
\[
\al_{\fh}(\la_p(X_p))=\la_p(X_p)
\foral p\in P\text{ and }\fh\in\fH.
\]
We write $e_P$ for the unit in $P\subseteq G$ and $e_\fH$ for the unit in $\fH$.

For each $p\in P$ define
\[
X_p\rtimes_{\al,\la}\fH
:=[C_c(\fH,\la_p(X_p))]
\subseteq \T_\la(X)\rtimes_{\al,\la}\fH.
\]
The family
\[
X\rtimes_{\al,\la}\fH:=\{X_p\rtimes_{\al,\la}\fH\}_{p\in P}
\]
defines a product system over $P$, with the reduced crossed product $A\rtimes_{\al,\la}\fH$ as its coefficient algebra; see for example
\cite[Section~3]{Kat20} and \cite[Lemma~7.11]{KR19}.  Let
\[
\io^\rtimes \colon X\rtimes_{\al,\la}\fH
\to \T_\la(X) \rtimes_{\al,\la}\fH
\]
denote the identity representation, and let $\wh{t}^\rtimes$ denote the universal
representation of $X\rtimes_{\al,\la}\fH$.

Let $\bo{x}$ be a constructible ideal in $\J$. Note that $\al_{\fh}(\la_p(X_p))=\la_p(X_p)$ for every $p\in P$ and hence $\al_{\fh}(\bo{K}_{\bo{x},\la_\ast}) =\bo{K}_{\bo{x},\la_\ast}$ for every $\fh \in \fH$. We set
\[
\bo{K}_{\bo{x},\la_\ast}\rtimes_{\al,\la}\fH:=[C_c(\fH,\bo{K}_{\bo{x},\la_\ast})] \subseteq \T_\la(X)\rtimes_{\al,\la}\fH.
\]

In order to make a distinction between the $\bo{K}$-cores of a representation of $X$ and the $\bo{K}$-cores of a representation of $X\rtimes_{\al,\la}\fH$, for a representation $t^\rtimes$ of $X\rtimes_{\al,\la}\fH$ we write $\bo{K}^\rtimes_{\bo{x},t^\rtimes_\ast}$ for the $\bo{K}$-core of the representation $t^\rtimes$ for a constructible ideal $\bo{x}\in \J$.

\begin{lemma}\label{L:kcores}
Let $\bo{x}$ be in $\J$. With the aforementioned notation, we have
\[
\bo{K}^\rtimes_{\bo{x},\io^\rtimes_\ast} \subseteq \bo{K}_{\bo{x},\la_\ast}\rtimes_{\al,\la}\fH.
\]
\end{lemma}

\begin{proof}
It suffices to consider an element $c\in \bo{K}^\rtimes_{\bo{x},\io^\rtimes_\ast}$ such that
\[
c=\io^\rtimes_{p_1}(f_{p_1})^* \io^\rtimes_{q_1}(h_{q_1})\cdots \io^\rtimes_{p_n}(f_{p_n})^* \io^\rtimes_{q_n}(h_{q_n})=f_{p_1}^* \ast h_{q_1} \ast \cdots \ast f_{p_n}^* \ast h_{q_n},
\]
where $n\in \bN$ and $p_1, q_1, \dots, p_n, q_n \in P$ satisfy
\[
p_1^{-1} q_1 \cdots p_n^{-1} q_n = e_P
\text{ and }
q_n^{-1} p_n \dots q_1^{-1} p_1 P = \bo{x},
\]
and $f_{p_i} \in C_c(\fH,\la_{p_i}(X_{p_i}))$ and $h_{q_i}\in C_c(\fH,\la_{q_i}(X_{q_i}))$. Since each $\la_p(X_p)$ is $\al$-invariant we get
\[
(f_{p_1}^* \ast h_{q_1})(\fg')= \int_{\fH} \De(\fg)^{-1} \al_\fg(f_{p_1}(\fg^{-1})^*) \al_\fg(h_{q_1}(\fg^{-1} \fg')) \, d\mu(\fg) \in [\la_{p_1}(X_{p_1})^* \la_{q_1}(X_{q_1})].
\]
Inductively, we obtain
\[
c\in C_c(\fH,[\la_{p_1}(X_{p_1})^* \la_{q_1}(X_{q_1})\cdots \la_{p_n}(X_{p_n})^* \la_{q_n}(X_{q_n})])\subseteq C_c(\fH, \bo{K}_{\bo{x},\la_\ast}) \subseteq \bo{K}_{\bo{x},\la_\ast}\rtimes_{\al,\la}\fH.
\]
The linear span of elements of this form is dense in $ \bo{K}^\rtimes_{\bo{x},\io^\rtimes_\ast}$ and hence the proof is complete.
\end{proof}

For the remainder of this section we fix a non-degenerate product system $X$ and a generalised gauge action $\al$ by a locally compact Hausdorff group. To ease notation let us write
\[
j:=i^r_{\T_\la(X)} \colon \T_\la(X) \to \M(\T_\la(X)\rtimes_{\al,\la}\fH)
\]
for the canonical embedding. Let us also write $\de^{\rtimes}$ for the reduced coaction by $\fH$ on $\T_\la(X)\rtimes_{\al,\la}\fH$, and
\[
\E_{\de^{\rtimes}}\colon \M_{\de^{\rtimes}} \to \M(\T_\la(X)\rtimes_{\al,\la}\fH)
\]
for the averaging map with respect to $\de^\rtimes$.

Since $X$ is non-degenerate, if $(e_i)_i$ is a contractive approximate unit for $A$, then $\la_{e_P}(e_i)$ is a contractive approximate unit for $\T_\la(X)$. In particular, by \cite[Lemma~3.5]{KR19} we may pick appropriate functions $f_i \in C_c(\fH)$ such that the functions $u_i\in C_c(\fH,\la_{e_P}(A))$ defined as 
\begin{equation}\label{Eq:appunit}
u_i(\fh)=f_i(\fh)\la_{e_P}(e_i) \foral \fh\in \fH,
\end{equation}
form a contractive approximate unit for $\T_\la(X)\rtimes_{\al,\la}\fH$.

We will need the following lemma.

\begin{lemma}\label{L:local1}
With the aforementioned notation, for each fixed $u_i$ defined in \eqref{Eq:appunit} we have:
\begin{enumerate}
\item $u_i \in \N_{\de^\rtimes}$.
\item If $z\in\T_\la(X)\rtimes_{\al,\la}\fH$, then $z u_i \in \N_{\de^\rtimes}$.

\item If $z,w\in \T_\la(X)\rtimes_{\al,\la}\fH$, then
$(z u_i)^*(w u_i)\in\M_{\de^{\rtimes}}$ and
$\E_{\de^{\rtimes}}((z u_i)^*(w u_i))\in j(\T_\la(X))$.
In particular, the map
\[
\Theta_i \colon
(\T_\la(X)\rtimes_{\al,\la}\fH)\times
(\T_\la(X)\rtimes_{\al,\la}\fH)\to j(\T_\la(X))
\text{ such that }
\Theta_i(z,w)
:=
\E_{\de^{\rtimes}}((zu_i)^*(w u_i)),
\]
is sesquilinear and
$\|\Theta_i(z,w)\|
\leq
\|z\|\cdot \|w\|\cdot \|\E_{\de^\rtimes}(u_i^*u_i)\|$.

\end{enumerate}
\end{lemma}

\begin{proof}
Item (i) follows from \eqref{Eq:Buss} since $u_i \in C_c(\fH,\la_{e_P}(A)) \subseteq C_c(\fH,\T_\la(X))$,
and item (ii) follows from the fact that 
$\N_{\de^{\rtimes}}$ is a left ideal in $\M(\T_\la(X)\rtimes_{\al,\la}\fH)$.

We now prove item (iii). 
Fix $z,w\in\T_\la(X)\rtimes_{\al,\la}\fH$. Since
$\M_{\de^{\rtimes}}=\spn\N_{\de^{\rtimes}}^*\N_{\de^{\rtimes}}$, item (ii) implies that $(z u_i)^*(w u_i)\in\M_{\de^{\rtimes}}$. Linearity of $\E_{\de^{\rtimes}}$ on $\M_{\de^{\rtimes}}$ shows that $\Theta_i$ is
sesquilinear.
Note that 
\[
0\leq u_i^* z^* z u_i
\leq \|z\|^2 u_i^* u_i,
\]
and therefore, since $\M_{\de^{\rtimes}}^+$ is a hereditary cone and $\E_{\de^{\rtimes}}$ is positive on $\M_{\de^\rtimes}^+$, we get
\[
0\leq
\E_{\de^{\rtimes}}((zu_i)^*(zu_i))
\leq
\|z\|^2 \E_{\de^\rtimes}(u_i^*u_i).
\]
Applying the Cauchy--Schwarz inequality of
\cite[Proposition~3.15]{KV99} to
\(\de^{\rtimes}(zu_i)\) and \(\de^{\rtimes}(wu_i)\) yields
\[
\begin{aligned}
\|\E_{\de^{\rtimes}}((zu_i)^*(wu_i))\|^2
&\leq
\|\E_{\de^{\rtimes}}((zu_i)^*(zu_i))\|\cdot
\|\E_{\de^{\rtimes}}((wu_i)^*(wu_i))\|\\
&\leq
\|z\|^2 \cdot \|w\|^2 \cdot \|\E_{\de^\rtimes}(u_i^*u_i)\|^2.
\end{aligned}
\]

It remains to prove that the range of $\Theta_i$ is a subset of $j(\T_\la(X))$. If
$z,w\in C_c(\fH,\T_\la(X))$, then 
\[
z \ast u_i,w \ast u_i\in C_c(\fH,\T_\la(X)),
\]
and hence \eqref{Eq:Buss} and \eqref{Eq:avercomp} give
\[
\Theta_i(z,w)
=
\E_{\de^{\rtimes}}((z \ast u_i)^* \ast (w \ast u_i))
\in j(\T_\la(X)).
\]
For arbitrary $z,w\in\T_\la(X)\rtimes_{\al,\la}\fH$ choose nets
$z_\nu,w_{\nu'}\in C_c(\fH,\T_\la(X))$ with $z_\nu \xrightarrow{\nu}z$ and $w_{\nu'} \xrightarrow{\nu'} w$ in norm. Note that $\|w_{\nu'}\| \to \|w\|$ and thus
\[
\|\Theta_i(z_\nu,w_{\nu'})-\Theta_i(z,w)\|
\leq
\|z_\nu-z\|\cdot \|w_{\nu'}\|\cdot \|\E_{\de^\rtimes}(u_i^*u_i)\|
+\|z\|\cdot \|w_{\nu'}-w\|\cdot \|\E_{\de^\rtimes}(u_i^*u_i)\| \xrightarrow{\nu,\nu'} 0,
\]
which completes the proof.
\end{proof}

Recall that for constructible ideals $\bo{x}, \bo{y}\in \J$ we have $[\bo{K}_{\bo{x},\la_\ast} \cdot \bo{K}_{\bo{y},\la_\ast}]\subseteq \bo{K}_{\bo{x}\cap \bo{y},\la_\ast}$.

\begin{lemma}\label{L:invar}
Let $\bo{x}, \bo{y} \in \J$ be constructible ideals of $P$. With the aforementioned notation, for each fixed $u_i$ defined in \eqref{Eq:appunit} we have that if $z\in \bo{K}_{\bo{x},\la_\ast}\rtimes_{\al,\la}\fH$ and
$w\in \bo{K}_{\bo{y},\la_\ast}\rtimes_{\al,\la}\fH$, then
$\E_{\de^{\rtimes}}((z u_i)^*(w u_i))\in j(\bo{K}_{\bo{x}\cap \bo{y},\la_\ast})$.
\end{lemma}

\begin{proof}
First suppose that $z\in C_c(\fH,\bo{K}_{\bo{x},\la_\ast})$ and
$w\in C_c(\fH,\bo{K}_{\bo{y},\la_\ast})$. We have
\[
(z \ast u_i)(\fg)=\int_{\fH}z(\fh) f_i(\fh^{-1}\fg) \al_{\fh}(\la_{e_P}(e_i))\,d\mu(\fh)
\]
and
\[
(w \ast u_i)(\fg)=\int_{\fH}w(\fh)f_i(\fh^{-1}\fg) \al_{\fh}(\la_{e_P}(e_i))\,d\mu(\fh).
\]
Since $\al_{\fh}(\la_{e_P}(A))=\la_{e_P}(A)$ for every $\fh\in \fH$ and $\bo{K}_{\bo{x},\la_\ast} \cdot \la_{e_P}(A) \subseteq \bo{K}_{\bo{x},\la_\ast}$, and similarly for $\bo{K}_{\bo{y},\la_\ast}$, we get $z \ast u_i\in C_c(\fH,\bo{K}_{\bo{x},\la_\ast})$ and $w \ast u_i\in C_c(\fH,\bo{K}_{\bo{y},\la_\ast})$.
By \eqref{Eq:Buss} and \eqref{Eq:avercomp} we get
\[
\E_{\de^{\rtimes}}((z \ast u_i)^* \ast (w\ast u_i))
=
j\left(\int_{\fH}\al_{\fh^{-1}}
\big((z \ast u_i)(\fh)^*(w \ast u_i)(\fh)\big)\,d\mu(\fh)\right).
\]
Since the closed subspace $[\bo{K}_{\bo{x},\la_\ast} \cdot \bo{K}_{\bo{y},\la_\ast} ]$ is $\al$-invariant the integral belongs to
\[
[\bo{K}_{\bo{x},\la_\ast} \cdot \bo{K}_{\bo{y},\la_\ast} ] \subseteq \bo{K}_{\bo{x}\cap \bo{y},\la_\ast}.
\]
For arbitrary
$z\in \bo{K}_{\bo{x},\la_\ast} \rtimes_{\al,\la}\fH$ and
$w\in \bo{K}_{\bo{y},\la_\ast} \rtimes_{\al,\la}\fH$ choose nets
$z_\nu\in C_c(\fH,\bo{K}_{\bo{x},\la_\ast})$ and $w_{\nu'}\in C_c(\fH,\bo{K}_{\bo{y},\la_\ast})$ with $z_\nu \xrightarrow{\nu}z$ and $w_{\nu'} \xrightarrow{\nu'} w$ in norm. By Lemma~\ref{L:local1} (iii) we have
\[
\|\Theta_i(z_\nu,w_{\nu'})-\Theta_i(z,w)\|
\xrightarrow{\nu,\nu'} 0.
\]
In particular, since $\Theta_i(z_\nu,w_{\nu'})\in j(\bo{K}_{\bo{x}\cap \bo{y},\la_\ast})$ we conclude that
\[
\Theta_i(z,w)=\E_{\de^{\rtimes}}((z u_i)^*(w u_i))\in j(\bo{K}_{\bo{x}\cap \bo{y},\la_\ast}),
\]
as required.
\end{proof}

We will also need the following lemma.

\begin{lemma}\label{L:local2}
With the aforementioned notation, let $(u_i)_i$ be the contractive approximate unit for $\T_\la(X)\rtimes_{\al,\la}\fH$ defined in \eqref{Eq:appunit}. If $z\in\T_\la(X) \rtimes_{\al,\la}\fH$ and
$\E_{\de^{\rtimes}}((z u_i)^*(zu_i))=0$
for every $i$,
then $z=0$.
\end{lemma}

\begin{proof}
If $\E_{\de^{\rtimes}}((z u_i)^*(zu_i))=0$ for all $i$, then we apply
Lemma~\ref{L:Bim} (ii) to get $z u_i=0$ for all $i$.
Since $z u_i \to z$ in norm, we get that $z=0$, as required.
\end{proof}


The following is the key result of this work.

\begin{theorem}\label{T:idisfock}
Let $X$ be a non-degenerate product system over a unital subsemigroup $P$ of a discrete group
$G$, and let $\al$ be a generalised gauge action of a locally compact Hausdorff group
$\fH$ on $\T_\la(X)$.  The identity representation
\[
\io^\rtimes \colon X\rtimes_{\al,\la}\fH
\to
\T_\la(X)\rtimes_{\al,\la}\fH
\]
is an injective Fock covariant representation of $X\rtimes_{\al,\la}\fH$ that admits a normal coaction by $G$.
\end{theorem}

\begin{proof}
Injectivity is immediate.  By \cite[Proposition 4.1]{DKKLL22} we have that $\T_\la(X)$ admits a normal coaction by $G$. Since $\al$ is a generalised gauge action there is an induced $*$-homomorphism
\[
\T_\la(X)\rtimes_{\al,\la}\fH \to (\T_\la(X) \otimes \ca(G)) \rtimes_{\al\otimes \id, \la} \fH.
\]
By \cite[Lemma~7.16]{Wil07} we have a canonical $*$-isomorphism
\[
 (\T_\la(X) \otimes \ca(G)) \rtimes_{\al\otimes \id, \la} \fH
 \simeq 
 (\T_\la(X) \rtimes_{\al, \la} \fH) \otimes \ca(G),
\]
and hence by composing we get that $\io^\rtimes$ admits a coaction by $G$. Normality of the coaction by $G$ on $\T_\la(X)$ yields in particular that $\io^\rtimes$ admits a normal coaction by $G$.

We will use the characterisation of Fock covariance of Theorem \ref{T:Fockcov}, and verify that $\io^\rtimes$ satisfies conditions (i) and (ii) therein. By \cite[Proposition 2.4]{KP25} we have $\bo{K}_{\mt,\la_\ast}=(0)$ and thus Lemma~\ref{L:kcores} yields that
\[
\bo{K}^\rtimes_{\mt,\io^\rtimes_\ast}
\subseteq
\bo{K}_{\mt,\la_\ast}\rtimes_{\al,\la}\fH=(0).
\]
This shows that condition (i) is satisfied by $\io^\rtimes$.

We now prove that condition (ii) is satisfied by $\io^\rtimes$.  Let
$\F=\{\bo{x}_1,\dots,\bo{x}_n\}\subseteq\J$ be finite and $\cap$-closed with
$\bigcup_{k=1}^n\bo{x}_k\neq\mt$, and
$b_k\in\bo{K}^{\rtimes}_{\bo{x}_k,\widehat t^\rtimes_\ast}$ such that
\[
\sum_{k:r\in\bo{x}_k}\io^\rtimes_\ast(b_k)\,
\io^\rtimes_r(X_r\rtimes_{\al,\la}\fH)=(0)
\foral r\in\bigcup_{k=1}^n\bo{x}_k.
\]
We must prove that $\sum_{k=1}^n\io^\rtimes_\ast(b_k)=0$.  To ease notation, set
\[
c:=\sum_{k=1}^n\io^\rtimes_\ast(b_k) \qand c_r:=\sum_{k:r\in\bo{x}_k}\io^\rtimes_\ast(b_k).
\]

Recall the contractive approximate unit $(u_i)_i \subseteq C_c(\fH,\la_{e_P}(A))$ for $\T_\la(X)\rtimes_{\al,\la}\fH$ defined in \eqref{Eq:appunit} and fix a $u_i$. By Lemma~\ref{L:kcores} we have that 
\[
\io^\rtimes_\ast(b_k) \in \bo{K}^{\rtimes}_{\bo{x}_k,\io^\rtimes_\ast} \subseteq \bo{K}_{\bo{x}_k,\la_\ast}\rtimes_{\al,\la}\fH \foral k\in \{1,\dots, n\}.
\]
Applying Lemma~\ref{L:invar} yields that
\[
\E_{\de^{\rtimes}}((\io^\rtimes_\ast(b_k) u_i)^*(\io^\rtimes_\ast(b_\ell) u_i))
\in j(\bo{K}_{\bo{x}_k\cap\bo{x}_\ell,\la_\ast}) \foral k,\ell\in \{1,\dots, n\}.
\]
Now choose $d_{k\ell}^{i}\in\bo{K}_{\bo{x}_k\cap\bo{x}_\ell,\widehat t_\ast}$ such that
\[
\E_{\de^{\rtimes}}((\io^\rtimes_\ast(b_k) u_i)^*(\io^\rtimes_\ast(b_\ell) u_i))
=j(\la_\ast(d_{k\ell}^{i})) \foral k,\ell\in \{1,\dots, n\}.
\]
By Lemma \ref{L:properties} we get that 
\[
\sum_{k,\ell=1}^n(\io^\rtimes_\ast(b_k) u_i)^*(\io^\rtimes_\ast(b_\ell) u_i) \in \M_{\de^{\rtimes}},
\]
and hence linearity of $\E_{\de^{\rtimes}}$ on $\M_{\de^{\rtimes}}$ gives us that
\[
\E_{\de^{\rtimes}}((c u_i)^*(c u_i))
=
\sum_{k,\ell=1}^n \E_{\de^{\rtimes}}((\io^\rtimes_\ast(b_k) u_i)^*(\io^\rtimes_\ast(b_\ell) u_i))
=\sum_{k,\ell=1}^n j(\la_\ast(d_{k\ell}^{i})).
\]
Moreover, for any $r\in\bigcup_{k=1}^n\bo{x}_k$ we have
\[
\{(k,\ell)\in \{1,\dots, n\}^2\mid r\in\bo{x}_k\cap\bo{x}_\ell\}
=
\{k\in \{1,\dots, n\} \mid r\in \bo{x}_k\} \times \{\ell \in \{1,\dots, n\} \mid r\in \bo{x}_\ell\},
\]
and hence
\[
\E_{\de^{\rtimes}}((c_ru_i)^*(c_ru_i))=\sum_{k:r\in \bo{x}_k} \sum_{\ell: r\in \bo{x}_\ell} \E_{\de^{\rtimes}}((\io^\rtimes_\ast(b_k) u_i)^*(\io^\rtimes_\ast(b_\ell) u_i))
= j\left(\sum_{k,\ell :r\in \bo{x}_k \cap \bo{x}_\ell} \la_\ast(d_{k\ell}^{i})\right).
\]

Now fix $r\in\bigcup_{k=1}^n\bo{x}_k$ and $\xi_r\in X_r$. We have $u_i j(\la_r(\xi_r))\in X_r\rtimes_{\al,\la}\fH$ (see for example \cite[Appendix A, Remark A.8]{EKQR06}) and thus by our hypothesis we get
\[
c_ru_i j(\la_r(\xi_r))=\sum_{k:r\in\bo{x}_k}\io^\rtimes_\ast(b_k)\,
\io^\rtimes_r(u_i j(\la_r(\xi_r)))=0.
\]
Applying Lemma~\ref{L:Bim} (i) for $(c_r u_i)^* (c_r u_i) \in \M_{\de^{\rtimes}}$ and $j(\la_r(\xi_r))\in \M^{\rm fix}_{\de^{\rtimes}}$ implies that
\[
0
=\E_{\de^{\rtimes}} (j(\la_r(\xi_r))^* (u_i^* c_r^* c_r u_i) j(\la_r(\xi_r)))
=j(\la_r(\xi_r))^* \E_{\de^{\rtimes}} ((c_ru_i)^*(c_ru_i))j(\la_r(\xi_r)).
\]
Note that $\E_{\de^{\rtimes}}((c_ru_i)^*(c_ru_i))\geq0$ and hence it follows that
\[
\E_{\de^{\rtimes}}((c_ru_i)^*(c_ru_i))^{1/2}j(\la_r(\xi_r))=0.
\]
In particular, since $\xi_r$ was arbitrary we get
\[
\E_{\de^{\rtimes}}((c_ru_i)^*(c_ru_i))j(\la_r(X_r))=0.
\]
Injectivity of $j$ then yields
\[
\sum_{k,\ell :r\in\bo{x}_k\cap\bo{x}_\ell}
\la_\ast(d_{k\ell}^{i})\la_r(X_r)=(0)
\foral r\in\bigcup_{k=1}^n\bo{x}_k=\bigcup_{k,\ell=1}^n\bo{x}_k\cap \bo{x}_\ell.
\]
Since $\la$ satisfies condition (ii) of
Theorem \ref{T:Fockcov} we get
\[
\sum_{k,\ell=1}^n\la_\ast(d_{k\ell}^{i})=0,
\]
and hence
\[
\E_{\de^{\rtimes}}((c u_i)^*(c u_i))=j\left(\sum_{k,\ell=1}^n\la_\ast(d_{k\ell}^{i})\right)=0.
\]
Since $i$ was arbitrary, Lemma~\ref{L:local2} then implies that $c=0$, as required.
\end{proof}

We are now ready to obtain the positive resolution of the reduced Hao--Ng isomorphism problem for non-degenerate product systems in full generality.

\begin{theorem}\label{T:HaoNg}
Let $P$ be a unital subsemigroup of a discrete group $G$ and let $X$ be a non-degenerate product system over $P$.
Let $\al$ be a generalised gauge action of a locally compact Hausdorff group $\fH$ on $\T_\la(X)$ and let $\dot{\al}$ denote the induced action of $\fH$ on $A \times_{X,\la} P$.
Then there exists a canonical completely isometric isomorphism
\[
\T_\la(X \rtimes_{\al, \la} \fH)^+ \simeq \T_\la(X)^+ \rtimes_{\al, \la} \fH,
\]
and consequently, a canonical $*$-isomorphism
\[
(A \rtimes_{\al, \la} \fH) \times_{X \rtimes_{\al, \la} \fH,\la} P
\simeq
(A \times_{X, \la} P) \rtimes_{\dot \al, \la} \fH.
\]
\end{theorem}

\begin{proof}
By Theorem~\ref{T:idisfock}, the identity representation
\[
\io^\rtimes:X\rtimes_{\al,\la}\fH
\to \T_\la(X)\rtimes_{\al,\la}\fH
\]
is injective and Fock covariant, and admits a normal coaction by $G$.
By \cite[Corollary 4.1]{KP25} we obtain a canonical completely isometric isomorphism between the tensor algebra $\T_\la(X\rtimes_{\al,\la}\fH)^+$ and the norm-closed algebra generated by $\{\io^\rtimes_p(X_p\rtimes_{\al,\la}\fH):p\in P\}$ inside $\T_\la(X)\rtimes_{\al,\la}\fH$. By \cite[Corollary 3.16]{KR19} we have that $\T_\la(X)^+ \rtimes_{\al, \la} \fH$ can be considered canonically inside $\T_\la(X) \rtimes_{\al, \la} \fH$, and it is exactly the norm-closed algebra generated by $\{\io^\rtimes_p(X_p\rtimes_{\al,\la}\fH):p\in P\}$. Therefore, we get a canonical completely isometric isomorphism
\[
\T_\la(X\rtimes_{\al,\la}\fH)^+
\simeq
\T_\la(X)^+ \rtimes_{\al, \la} \fH.
\]
Since $X$ is non-degenerate, any contractive approximate unit for $A$ gives a 
contractive approximate unit for $\T_\la(X)^+$.
By \cite[Theorem~4.4]{DT26} we get a canonical
$*$-isomorphism
\[
\cenv(\T_\la(X)^+\rtimes_{\al,\la}\fH)
\simeq
\cenv(\T_\la(X)^+)\rtimes_{\dot\al,\la}\fH,
\]
and \cite[Theorem~5.1]{Seh22} completes the proof.
\end{proof}


\begin{thebibliography}{99}

\bibitem{Aba10}
B.\ Abadie,
\textit{Takai duality for crossed products by Hilbert $C^*$-bimodules},
J.\ Operator Theory \textbf{64} (2010), no.\ 1, 19--34.

\bibitem{Arv89}
W.\ B.\ Arveson,
\textit{Continuous analogues of Fock space},
Mem.\ Amer.\ Math.\ Soc.\ \textbf{80} (1989), no.\ 409, iv+66 pp.

\bibitem{Arv11}
W.\ B.\ Arveson,
\textit{The noncommutative Choquet boundary II: hyperrigidity},
Israel J.\ Math.\ \textbf{184} (2011), 349--385.

\bibitem{BKQR15}
E.\ B\'edos, S.\ Kaliszewski, J.\ Quigg and D.\ Robertson,
\textit{A new look at crossed product correspondences and associated
$C^*$-algebras},
J.\ Math.\ Anal.\ Appl.\ \textbf{426} (2015), no.\ 2, 1080--1098.

\bibitem{BL04}
D.\ P.\ Blecher and C.\ Le Merdy,
\textit{Operator algebras and their modules---an operator space approach},
London Mathematical Society Monographs, New Series, vol.\ 30,
Oxford University Press, Oxford, 2004.

\bibitem{BLS18}
N.\ Brownlowe, N.\ S.\ Larsen and N.\ Stammeier,
\textit{$C^*$-algebras of algebraic dynamical systems and right LCM semigroups},
Indiana Univ.\ Math.\ J.\ \textbf{67} (2018), no.\ 6, 2453--2486.

\bibitem{BM09}
A.\ Buss and R.\ Meyer,
\textit{Square-integrable coactions of locally compact quantum groups},
Rep.\ Math.\ Phys.\ \textbf{63} (2009), no.\ 1, 191--224.

\bibitem{Bus10}
A.\ Buss,
\textit{Integrability of dual coactions on Fell bundle $C^*$-algebras},
Bull.\ Braz.\ Math.\ Soc.\ (N.S.) \textbf{41} (2010), no.\ 4, 607--641.

\bibitem{CLSV11}
T.\ M.\ Carlsen, N.\ S.\ Larsen, A.\ Sims and S.\ T.\ Vittadello,
\textit{Co-universal algebras associated to product systems, and
gauge-invariant uniqueness theorems},
Proc.\ Lond.\ Math.\ Soc.\ (3) \textbf{103} (2011), no.\ 4, 563--600.

\bibitem{Com68}
F.\ Combes,
\textit{Poids sur une $C^*$-alg\`ebre},
J.\ Math.\ Pures Appl.\ (9) \textbf{47} (1968), 57--100 (French).

\bibitem{CDL13}
J.\ Cuntz, C.\ Deninger and M.\ Laca,
\textit{$C^*$-algebras of Toeplitz type associated with algebraic number fields},
Math.\ Ann.\ \textbf{355} (2013), no.\ 4, 1383--1423.

\bibitem{Din91}
H.\ T.\ Dinh,
\textit{Discrete product systems and their $C^*$-algebras},
J.\ Funct.\ Anal.\ \textbf{102} (1991), no.\ 1, 1--34.

\bibitem{DKKLL22}
A.\ Dor-On, E.\ T.\ A.\ Kakariadis, E.\ G.\ Katsoulis, M.\ Laca and X.\ Li,
\textit{$C^*$-envelopes for operator algebras with a coaction and co-universal
$C^*$-algebras for product systems},
Adv.\ Math.\ \textbf{400} (2022), Paper No.\ 108286, 40 pp.

\bibitem{DK20}
A.\ Dor-On and E.\ G.\ Katsoulis,
\textit{Tensor algebras of product systems and their $C^*$-envelopes},
J.\ Funct.\ Anal.\ \textbf{278} (2020), no.\ 7, Paper No.\ 108416, 32 pp.

\bibitem{DT26}
A.\ Dor-On and I.\ Thompson,
\textit{The Hao--Ng isomorphism theorem for reduced crossed products},
preprint, arXiv:2505.00587v3 (2026).

\bibitem{EKQR06}
S.\ Echterhoff, S.\ Kaliszewski, J.\ Quigg and I.\ Raeburn,
\textit{A categorical approach to imprimitivity theorems for $C^*$-dynamical systems},
Mem.\ Amer.\ Math.\ Soc.\ \textbf{180} (2006), no.\ 850, viii+169 pp.

\bibitem{EN02}
R.\ Exel and C.-K.\ Ng,
\textit{Approximation property of $C^*$-algebraic bundles},
Math.\ Proc.\ Cambridge Philos.\ Soc.\ \textbf{132} (2002), no.\ 3, 509--522.

\bibitem{Exe17}
R.\ Exel,
\textit{Partial dynamical systems, Fell bundles and applications},
Mathematical Surveys and Monographs, vol.\ 224,
American Mathematical Society, Providence, RI (2017) vi+321 pp.

\bibitem{FD88}
J.\ M.\ G.\ Fell and R.\ S.\ Doran,
\textit{Representations of $*$-algebras, locally compact groups, and Banach
$*$-algebraic bundles. Vol.\ 2},
Pure and Applied Mathematics, vol.\ 126,
Academic Press, Boston, MA, 1988.

\bibitem{Fol16}
G.\ B.\ Folland,
\textit{A course in abstract harmonic analysis},
2nd ed.,
CRC Press, Boca Raton, FL, 2016.

\bibitem{Fow02}
N.\ J.\ Fowler,
\textit{Discrete product systems of Hilbert bimodules},
Pacific J.\ Math.\ \textbf{204} (2002), no.\ 2, 335--375.

\bibitem{FR98}
N.\ J.\ Fowler and I.\ Raeburn,
\textit{Discrete product systems and twisted crossed products by semigroups},
J.\ Funct.\ Anal.\ \textbf{155} (1998), no.\ 1, 171--204.

\bibitem{Ham79}
M.\ Hamana,
\textit{Injective envelopes of operator systems},
Publ.\ Res.\ Inst.\ Math.\ Sci.\ \textbf{15} (1979), no.\ 3, 773--785.

\bibitem{HN08}
G.\ Hao and C.-K.\ Ng,
\textit{Crossed products of $C^*$-correspondences by amenable group actions},
J.\ Math.\ Anal.\ Appl.\ \textbf{345} (2008), no.\ 2, 702--707.

\bibitem{KKLL22}
E.\ T.\ A.\ Kakariadis, E.\ G.\ Katsoulis, M.\ Laca and X.\ Li,
\textit{Boundary quotient $C^*$-algebras of semigroups},
J.\ Lond.\ Math.\ Soc.\ (2) \textbf{105} (2022), no.\ 4, 2136--2166.

\bibitem{KKLL23} E.\ T.\ A.\ Kakariadis,  E.\ G.\ Katsoulis, M.\ Laca and X.\ Li,
\textit{Co-universality and controlled maps on product systems by right LCM-semigroups}, 
Analysis \& PDE \textbf{16} (2023), no.\ 6, 1433--1483.

\bibitem{KP25}
E.\ T.\ A.\ Kakariadis and I.\ A.\ Paraskevas,
\textit{On Fock covariance for product systems and the reduced Hao--Ng
isomorphism problem by discrete actions},
Proc.\ Roy.\ Soc.\ Edinburgh Sect.\ A, First View (2025), 1--53,
doi:10.1017/prm.2025.21.

\bibitem{KMQW10}
S.\ Kaliszewski, P.\ S.\ Muhly, J.\ Quigg and D.\ P.\ Williams,
\textit{Coactions and Fell bundles},
New York J.\ Math.\ \textbf{16} (2010), 315--359.

\bibitem{Kat84}
Y.\ Katayama,
\textit{Takesaki's duality for a non-degenerate co-action},
Math.\ Scand.\ \textbf{55} (1984), no.\ 1, 141--151.

\bibitem{Kat17}
E.\ G.\ Katsoulis,
\textit{$C^*$-envelopes and the Hao--Ng isomorphism for discrete groups},
Int.\ Math.\ Res.\ Not.\ IMRN \textbf{2017} (2017), no.\ 18, 5751--5768.

\bibitem{Kat20}
E.\ G.\ Katsoulis,
\textit{Product systems of $C^*$-correspondences and Takai duality},
Israel J.\ Math.\ \textbf{240} (2020), no.\ 1, 223--251.

\bibitem{KK06}
E.\ G.\ Katsoulis and D.\ W.\ Kribs,
\textit{Tensor algebras of $C^*$-correspondences and their $C^*$-envelopes},
J.\ Funct.\ Anal.\ \textbf{234} (2006), no.\ 1, 226--233.

\bibitem{KR19}
E.\ G.\ Katsoulis and C.\ Ramsey,
\textit{Crossed products of operator algebras},
Mem.\ Amer.\ Math.\ Soc.\ \textbf{258} (2019), no.\ 1240, vii+85 pp.

\bibitem{KR21}
E.\ G.\ Katsoulis and C.\ Ramsey,
\textit{The non-selfadjoint approach to the Hao--Ng isomorphism},
Int.\ Math.\ Res.\ Not.\ IMRN \textbf{2021} (2021), no.\ 2, 1160--1197.

\bibitem{Kat03}
T.\ Katsura,
\textit{A construction of $C^*$-algebras from $C^*$-correspondences},
in \textit{Advances in quantum dynamics}, 173--182,
Contemp.\ Math., vol.\ 335,
Amer.\ Math.\ Soc., Providence, RI (2003).

\bibitem{Kat04}
T.\ Katsura,
\textit{On $C^*$-algebras associated with $C^*$-correspondences},
J.\ Funct.\ Anal.\ \textbf{217} (2004), no.\ 2, 366--401.

\bibitem{KP99}
A.\ Kumjian and D.\ Pask,
\textit{$C^*$-algebras of directed graphs and group actions},
Ergodic Theory Dynam.\ Systems \textbf{19} (1999), no.\ 6, 1503--1519.

\bibitem{Kus97A}
J.\ Kustermans,
\textit{KMS-weights on $C^*$-algebras},
preprint, arXiv:funct-an/9704008 (1997).

\bibitem{KV99}
J.\ Kustermans and S.\ Vaes,
\textit{Weight theory for $C^*$-algebraic quantum groups},
preprint, arXiv:math/9901063 (1999).

\bibitem{KL19a}
B.\ K.\ Kwa\'{s}niewski and N.\ S.\ Larsen,
\textit{Nica--Toeplitz algebras associated with right-tensor $C^*$-precategories
over right LCM semigroups},
Internat.\ J.\ Math.\ \textbf{30} (2019), no.\ 2, 1950013, 57 pp.

\bibitem{KL19b}
B.\ K.\ Kwa\'{s}niewski and N.\ S.\ Larsen,
\textit{Nica--Toeplitz algebras associated with product systems over right
LCM semigroups},
J.\ Math.\ Anal.\ Appl.\ \textbf{470} (2019), no.\ 1, 532--570.

\bibitem{LS22}
M.\ Laca and C.\ F.\ Sehnem,
\textit{Toeplitz algebras of semigroups},
Trans.\ Amer.\ Math.\ Soc.\ \textbf{375} (2022), no.\ 10, 7443--7507.

\bibitem{Lan79}
M.\ B.\ Landstad,
\textit{Duality theory for covariant systems},
Trans.\ Amer.\ Math.\ Soc.\ \textbf{248} (1979), 223--267.

\bibitem{Lan95}
E.\ C.\ Lance,
\textit{Hilbert $C^*$-modules: a toolkit for operator algebraists},
London Mathematical Society Lecture Note Series, vol.\ 210,
Cambridge University Press, Cambridge, 1995.

\bibitem{Li12}
X.\ Li,
\textit{Semigroup $C^*$-algebras and amenability of semigroups},
J.\ Funct.\ Anal.\ \textbf{262} (2012), no.\ 10, 4302--4340.

\bibitem{Ng96}
C.-K.\ Ng,
\textit{Discrete coactions on $C^*$-algebras},
J.\ Austral.\ Math.\ Soc.\ Ser.\ A \textbf{60} (1996), no.\ 1, 118--127.

\bibitem{Nic92}
A.\ Nica,
\textit{$C^*$-algebras generated by isometries and Wiener--Hopf operators},
J.\ Operator Theory \textbf{27} (1992), no.\ 1, 17--52.

\bibitem{Nor14}
M.\ D.\ Norling,
\textit{Inverse semigroup $C^*$-algebras associated with left cancellative semigroups},
Proc.\ Edinb.\ Math.\ Soc.\ (2) \textbf{57} (2014), no.\ 2, 533--564.

\bibitem{Pau02}
V.\ I.\ Paulsen,
\textit{Completely bounded maps and operator algebras},
Cambridge Studies in Advanced Mathematics, vol.\ 78, xii+300 pp.,
Cambridge University Press, Cambridge (2002).

\bibitem{Ped79}
G.\ K.\ Pedersen,
\textit{$C^*$-algebras and their automorphism groups},
London Mathematical Society Monographs, vol.\ 14,
Academic Press, London (1979).

\bibitem{Pim97}
M.\ V.\ Pimsner,
\textit{A class of $C^*$-algebras generalizing both Cuntz--Krieger algebras and
crossed products by $\mathbb{Z}$},
in \textit{Free probability theory} (Waterloo, ON, 1995), 189--212,
Fields Inst.\ Commun., vol.\ 12,
Amer.\ Math.\ Soc., Providence, RI (1997).

\bibitem{Seh19}
C.\ F.\ Sehnem,
\textit{On $C^*$-algebras associated to product systems},
J.\ Funct.\ Anal.\ \textbf{277} (2019), no.\ 2, 558--593.

\bibitem{Seh22}
C.\ F.\ Sehnem,
\textit{$C^*$-envelopes of tensor algebras of product systems},
J.\ Funct.\ Anal.\ \textbf{283} (2022), no.\ 12, Paper No.\ 109707.

\bibitem{Tak03}
M.\ Takesaki,
\textit{Theory of operator algebras II},
Encyclopaedia of Mathematical Sciences, vol.\ 125,
Springer, Berlin (2003).

\bibitem{Wil07}
D.\ P.\ Williams,
\textit{Crossed products of $C^*$-algebras},
Mathematical Surveys and Monographs, vol.\ 134,
American Mathematical Society, Providence, RI (2007).

\end{thebibliography}

\end{document}
